\documentclass[11pt]{article}
\usepackage[margin=1in]{geometry}
\usepackage{hyperref}
\usepackage{graphicx}
\graphicspath{{./}}
\usepackage{natbib}
\usepackage{doi}
% \setcitestyle{numbers,super} % Superscript citation numbers (e.g. ¹, ²⁻⁴) — does not support locator notes
\setcitestyle{numbers,square} % Square brackets around citation numbers (e.g. [1], [2-4]) — math standard
% ── Merge citation text into a single hyperlink ───────────────────────────────
\makeatletter
\renewcommand{\hyper@natlinkbreak}[2]{#1}
\makeatother
\usepackage{titlesec}
% New page before each top-level section
% \newcommand{\sectionbreak}{\clearpage}

% ── Theorem-like environments ─────────────────────────────────────────────────
\usepackage{amsmath}
\usepackage{amsthm}
\usepackage{dsfont} % double-struck numerals (\mathds{1} for indicator functions)
% thmtools is only needed for the optional \listoftheorems feature. Loading it
% patches \newtheorem, and thmtools 0.76's patch is incompatible with recent
% LaTeX kernels (>= 2025) for shared-counter theorems (every sibling raises
% "Command \c@... already defined"). We therefore load it only when the
% list-of-theorems feature is actually requested.
% Theorem-like environments defined with amsthm's native shared-counter syntax
% (\newtheorem{child}[parent]{Name}). Numbering is grouped by counter without
% relying on thmtools' sibling mechanism.
\theoremstyle{plain}
\newtheorem{theorem}{Theorem}
\newtheorem{proposition}[theorem]{Proposition}
\newtheorem{lemma}[theorem]{Lemma}
\newtheorem{corollary}[theorem]{Corollary}
\newtheorem{conjecture}[theorem]{Conjecture}

\theoremstyle{definition}
\newtheorem{definition}{Definition}
\newtheorem{axiom}[definition]{Axiom}
\newtheorem{assumption}[definition]{Assumption}
\newtheorem{criterion}[definition]{Criterion}
\newtheorem{example}[definition]{Example}

\theoremstyle{remark}
\newtheorem{remark}{Remark}
\newtheorem{observation}[remark]{Observation}
\newtheorem{property}[remark]{Property}

% Prevent theorem-like environments from starting near the page bottom.
\usepackage{etoolbox}
\BeforeBeginEnvironment{theorem}{\needspace{4\baselineskip}\addvspace{6pt}}
\BeforeBeginEnvironment{corollary}{\needspace{4\baselineskip}\addvspace{6pt}}
\BeforeBeginEnvironment{lemma}{\needspace{4\baselineskip}\addvspace{6pt}}
\BeforeBeginEnvironment{proposition}{\needspace{4\baselineskip}\addvspace{6pt}}
\BeforeBeginEnvironment{definition}{\needspace{4\baselineskip}\addvspace{6pt}}
\BeforeBeginEnvironment{axiom}{\needspace{4\baselineskip}\addvspace{6pt}}
\BeforeBeginEnvironment{assumption}{\needspace{4\baselineskip}\addvspace{6pt}}
\BeforeBeginEnvironment{remark}{\needspace{4\baselineskip}\addvspace{9pt}}

% ── Pandoc-generated table support ────────────────────────────────────────────
\usepackage{longtable}
\usepackage{booktabs}
\usepackage{array}
\usepackage{calc}
\usepackage{caption}
% General table note: small font, flush-left, below the table (used by some papers).
\newenvironment{tablenote}
  {\par\addvspace{2pt}\footnotesize\noindent}
  {\par\addvspace{4pt}}
% Pandoc 3.x wraps uncaptioned longtables in \def\LTcaptype{none};
% this requires a 'none' counter to exist.
\newcounter{none}

% ── Paragraph spacing (matches Pandoc default) ────────────────────────────────
% No paragraph indentation; inter-paragraph spacing instead.
\usepackage{parskip}

% ── Page-break quality ────────────────────────────────────────────────────────
% Prevent widow/orphan lines (single line stranded at top/bottom of page).
\widowpenalty=10000
\clubpenalty=10000
% Pin footnotes to the page bottom even with \raggedbottom.
\usepackage[bottom]{footmisc}
% Prevent page breaks immediately before or after display equations.
\predisplaypenalty=10000
\postdisplaypenalty=150
% Allow uneven page bottoms rather than stretching vertical glue.
\raggedbottom
% Increase separation between body text and footnotes.
\setlength{\skip\footins}{1.5\baselineskip}
% Ensure headings don't appear at page bottom without body text following.
\usepackage{needspace}
% Reserve enough space for the heading itself plus at least two following
% lines of body text (or a nested heading plus a line). This prevents a
% heading from appearing at the bottom of a page with its content (or its
% next-level subheading) starting on the following page.
\let\origsection\section
\renewcommand{\section}{\needspace{6\baselineskip}\origsection}
\let\origsubsection\subsection
\renewcommand{\subsection}{\needspace{6\baselineskip}\origsubsection}
\let\origsubsubsection\subsubsection
\renewcommand{\subsubsection}{\needspace{4\baselineskip}\origsubsubsection}
\let\origparagraphcmd\paragraph
\renewcommand{\paragraph}{\needspace{3\baselineskip}\origparagraphcmd}
\let\origsubparagraph\subparagraph
\renewcommand{\subparagraph}{\needspace{3\baselineskip}\origsubparagraph}

% Discourage page breaks inside multi-line math environments (align, gather, etc.)
\interdisplaylinepenalty=10000

% ── Section numbering ────────────────────────────────────────────────────────
% Pandoc --number-sections generates \section{...} (not \section*{...}).
% LaTeX auto-numbers to depth 3 (subsubsection) by default.
\setcounter{secnumdepth}{3}

% ── Make \paragraph and \subparagraph block-level (matches Pandoc) ───────────
% By default LaTeX renders \paragraph as a run-in heading (text continues on
% the same line). These commands make them free-standing headings with
% proper spacing before and after.
\titleformat{\paragraph}
  {\normalfont\normalsize\bfseries}{\theparagraph}{1em}{}
\titlespacing*{\paragraph}
  {0pt}{3.25ex plus 1ex minus .2ex}{1.5ex plus .2ex}
\titleformat{\subparagraph}
  {\normalfont\normalsize\bfseries}{\thesubparagraph}{1em}{}
\titlespacing*{\subparagraph}
  {0pt}{3.25ex plus 1ex minus .2ex}{1.5ex plus .2ex}

% ── Pandoc-generated tightlist support ────────────────────────────────────────
\providecommand{\tightlist}{%
  \setlength{\itemsep}{0pt}\setlength{\parskip}{0pt}}

% ── Pandoc 3.x figure sizing ─────────────────────────────────────────────────
% \pandocbounded constrains included graphics to text width when they exceed it
\newcommand{\pandocbounded}[1]{%
  \sbox0{#1}%
  \ifdim\wd0>\linewidth
    \resizebox{\linewidth}{!}{#1}%
  \else
    #1%
  \fi
}

% ── Unicode character fixes ───────────────────────────────────────────────────
\usepackage{amssymb}
\usepackage{newunicodechar}
\newunicodechar{✓}{$\checkmark$}
\newunicodechar{≠}{$\neq$}
\newunicodechar{≡}{$\equiv$}
\newunicodechar{≈}{$\approx$}
\newunicodechar{δ}{$\delta$}
\newunicodechar{−}{$-$}
% Superscript digits used in scientific notation (e.g. 10^-7, 10^29)
\newunicodechar{⁰}{$^{0}$}
% ¹²³ already defined by textcomp; suppress redefinition warning
\begingroup\renewcommand\PackageWarning[2]{}%
\newunicodechar{¹}{$^{1}$}
\newunicodechar{²}{$^{2}$}
\newunicodechar{³}{$^{3}$}
\endgroup
\newunicodechar{⁴}{$^{4}$}
\newunicodechar{⁵}{$^{5}$}
\newunicodechar{⁶}{$^{6}$}
\newunicodechar{⁷}{$^{7}$}
\newunicodechar{⁸}{$^{8}$}
\newunicodechar{⁹}{$^{9}$}
\newunicodechar{⁻}{$^{-}$}
% Fix \rightarrow for use outside math mode
\let\mystRightarrow\rightarrow
\renewcommand{\rightarrow}{\ensuremath{\mystRightarrow}}
% ─────────────────────────────────────────────────────────────────────────────

% Allow TeX extra stretch to resolve moderate overfull hboxes
\emergencystretch=2em

\hypersetup{colorlinks=true, allcolors=blue}

\title{Thermodynamics of Cooperation: Value Dynamics}
\author{Keith Lostracco}
\date{}

\begin{document}
\maketitle

\begin{abstract}
Self-maintaining agents that share finite resources must couple to
high-energy centers (accumulated-negentropy nodes) for energy inflow,
yet tight coupling imposes autonomy costs that can erode the agent's
boundary integrity. We model this trade-off as a one-dimensional
potential on the coupling distance between an agent and a center,
combining a short-range-dominant dissolution cost (\(1/r^2\)) with a
long-range resource inflow (\(1/r\)) and a constant entropy maintenance
baseline. The resulting coexistence potential possesses a unique global
minimum, the coexistence attractor, whose global asymptotic stability we
prove via Lyapunov analysis. The zero-crossings of the potential define
the stable coexistence band, a bounded interval of coupling distances
within which the agent maintains positive net energy. We derive a
closed-form expression for the bandwidth of this band (the freedom
bandwidth) and prove that it increases with center mass and decreases
with agent complexity, dissolution pressure, and entropy leakage.
Outside the band, agents undergo irreversible transitions: dissolution
(boundary erosion from excessive coupling) or starvation (energy deficit
from insufficient coupling). We extend the single-center model to \(K\)
centers, proving that the multi-center potential separates across
coupling dimensions, yielding a componentwise global attractor with a
diversification benefit: agents that would be non-viable at any single
center may survive through distributed coupling.
\end{abstract}
\section{Introduction}\label{introduction}

The preceding papers in this series derive two classes of results.
Papers I through IV \citep{TC-I, TC-II, TC-III, TC-IV} prove that mutual
boundary constraints are necessary for any feasible allocation among
self-maintaining systems sharing finite resources, that violating these
constraints is thermodynamically self-destructive, that strategic
corruption of communication channels imposes super-linear efficiency
losses, and that complex systems represent irreplaceable accumulated
negentropy. \emph{Cooperative Equilibrium} \citep{TC-V} proves that
under energy-denominated payoffs with physical friction costs, universal
cooperation is the unique Nash Equilibrium that is simultaneously
Pareto-efficient and welfare-maximizing for agents with discount factor
\(\delta \geq 0.431\).

These results establish the static conditions under which cooperation is
optimal; they do not address the dynamical question of how agents reach
and maintain cooperative configurations. This paper fills that gap.

The principal contributions are:

\begin{enumerate}
\def\labelenumi{\arabic{enumi}.}
\item
  A \textbf{dynamical systems model} of agent-center coupling
  (§\ref{the-agent-center-model}), proving that stable coexistence is a
  global attractor (Theorem~\ref{thm-stability-coexistence-attractor})
  with a measurable freedom bandwidth
  (Theorem~\ref{thm-freedom-bandwidth}), and a coupled
  boundary-integrity dynamics with an explicit basin of no return
  (Theorem~\ref{thm-basin-of-no-return}) governing dissolution and
  starvation (Theorem~\ref{thm-irreversibility-dissolution},
  Theorem~\ref{thm-starvation-spiral}).
\item
  A \textbf{multi-center extension} (§\ref{multi-center-dynamics})
  proving separable optimization with diversification benefits
  (Theorem~\ref{thm-multi-center-attractor}), and a
  stability-cooperation feedback loop
  (Proposition~\ref{prop-stability-cooperation-feedback}) connecting the
  dynamical model to the game-theoretic equilibrium.
\item
  \textbf{Comparative statics} characterizing how freedom bandwidth
  depends on center mass, agent complexity, and coupling parameters
  (§\ref{comparative-statics-and-inequality}), with implications for
  inequality across heterogeneous agent populations.
\item
  \textbf{Irreversibility results} establishing a basin of no return for
  boundary integrity: the boundary is modeled as an active dissipative
  structure whose repair capacity and mobility scale with its own
  integrity, so below a critical integrity \(B_c\) both vanish together
  and the boundary reaches its collapse floor in finite time
  (§\ref{boundary-dynamics-and-irreversibility}). Dissolution and
  starvation are one-way processes once depletion crosses the basin
  boundary --- not at the band boundaries themselves --- and the repair
  side of the dynamics carries the repair multiplier \(\kappa > 1\) of
  \emph{Thermodynamic Friction} \citep{TC-IV}, grounding the
  irreversibility thermodynamically.
\end{enumerate}

\subsection{Related Work}\label{related-work}

\textbf{Dynamical systems and agent-center models.} The coexistence
potential introduced in §\ref{the-agent-center-model} draws on the
mathematical toolkit of gradient dynamics and Lyapunov stability
\citep{Strogatz2015, HirschSmaleDevaney2012}. The two-body potential
form (short-range repulsion from dissolution pressure, long-range
attraction from resource access) parallels the Lennard-Jones potential
in molecular physics \citep{Jones1924}, adapted here to socio-economic
coupling. The parallel is mathematical, not physical: the shared
structure is a potential with competing short-range and long-range
terms, not a claim that intermolecular forces operate at institutional
scales. Organizational economics has long modeled firm-worker
relationships as vertical coupling with autonomy costs and resource
benefits \citep{Tirole1988}, and institutional economics has documented
the concentration dynamics that the model formalizes
\citep{AcemogluRobinson2012}.

\textbf{Multi-agent economics and organizational theory.} The
agent-center model formalizes the vertical coupling between individuals
and resource-providing institutions. The comparative statics results
connect to classical questions in organizational economics about
autonomy versus integration \citep{Tirole1988}, and formalize the
destructive autonomy costs of absolutist, extractive regimes, while the
multi-center extension provides a dynamical foundation for the stability
and survival benefits of institutional pluralism
\citep{AcemogluRobinson2012}. The freedom bandwidth theorem quantifies a
trade-off familiar in labor economics: agents with higher complexity
(higher boundary maintenance costs) face narrower viable operating
ranges, a formal analog of the vulnerability of specialized workers to
market shocks.

\textbf{Mathematical ecology.} The coexistence band and its dependence
on center mass parallels the viable habitat ranges in population
ecology, where species occupy niches defined by resource availability
and metabolic cost \citep{Tilman1982}. The irreversibility results
(dissolution and starvation spirals) formalize the extinction vortex
phenomenon, in which small populations face self-reinforcing decline
\citep{Gilpin1986}. The multi-center diversification benefit provides a
mathematical foundation for the empirical observation that species with
broader resource bases are more resilient to perturbation.

\section{The Agent-Center Model}\label{the-agent-center-model}

\subsection{Motivation: Why Dynamics
Matter}\label{motivation-why-dynamics-matter}

The preceding derivations \citep{TC-I, TC-II, TC-III, TC-IV, TC-V}
established the static structure of cooperative multi-agent systems:
boundary constraints as necessary conditions for feasible allocation,
unilateral violation as thermodynamically wasteful, universal
cooperation as the unique efficient equilibrium, and information
integrity as a precondition for coordination. But the theory describes a
\textbf{dynamic} phenomenon: accumulated value creates attractor
dynamics. Entities are drawn toward high-energy centers because
proximity lowers the cost of fighting entropy. Too close, and the entity
is assimilated; too far, and it starves.

This is inherently a \textbf{dynamical systems} question: given an agent
and a center, what coupling strength is stable? What happens when the
agent is perturbed from equilibrium? Under what conditions does the
equilibrium exist at all?

\subsection{Primitive Objects}\label{primitive-objects}

\begin{definition}[High-Energy Center]
\label{def-high-energy-center}

A \emph{high-energy center} (HEC) is a network node with accumulated value mass

$$\mathcal{M} > 0 \quad \text{(energy units, J)}$$

representing the total accumulated negentropy of the center (Definition 6 \citep{TC-III}). The HEC may be a corporation, a state, a trade hub, or an ecosystem, i.e., any entity or coalition that has aggregated sufficient energy and information to exert significant influence on surrounding agents.
\end{definition}

\begin{definition}[Coupling Distance]
\label{def-coupling-distance}

The \emph{coupling distance} $r \in (0, \infty)$ between an agent $i$ and a HEC is a generalized coordinate measuring the degree of the agent's independence from the center:

\begin{itemize}
\item $r \to 0$: maximal coupling (complete integration, e.g., employee absorbed into corporation, citizen fully subject to state, subsidiary merged into parent)
\item $r \to \infty$: minimal coupling (complete isolation, e.g., autarkic entity, no trade or institutional affiliation)
\end{itemize}
The coupling distance is \textbf{not} a spatial distance. It is a scalar summarizing the composite engagement level across economic, informational, political, and social dimensions. Formally, $r$ can be constructed as a norm on the vector of coupling channels, but for the 1D analysis below, only its scalar properties matter.
\end{definition}

\textbf{Interpretation:} The coupling distance \(r\) is the single
degree of freedom in our model. An agent's ``strategy'' in the
value-dynamics context is its choice of how tightly to bind itself to a
given HEC. This complements the multi-dimensional resource strategy
\(\mathbf{x}_i\) from the Lagrangian constraints derivation: there, the
agent chose \emph{what} to consume; here, the agent chooses \emph{how
closely} to affiliate with the source of resources.

\subsection{The Resource Inflow
Function}\label{the-resource-inflow-function}

As the agent couples more tightly with a HEC (smaller \(r\)), it gains
greater access to the center's accumulated resources, including shared
infrastructure, trade networks, institutional protection, information
flows, and surplus energy. We model the \textbf{energy inflow rate} as:

\[E_{\text{in}}(r) = \frac{G \, \mathcal{M}}{r}\]

where \(G > 0\) is the \textbf{resource coupling coefficient}, a
proportionality constant capturing the efficiency of the HEC at
distributing resources to its affiliates.

\textbf{Properties:} - \(E_{\text{in}}(r) > 0\) for all \(r > 0\) (the
HEC always provides some resources, however distant). -
\(E_{\text{in}}(r) \to \infty\) as \(r \to 0\) (arbitrarily large
resource access at arbitrarily tight coupling, but offset by costs
below). - \(E_{\text{in}}(r) \to 0\) as \(r \to \infty\) (isolation
yields negligible resource access). - \(E_{\text{in}}\) is proportional
to \(\mathcal{M}\): a more massive center supplies more resources at
every coupling distance.

The \(1/r\) functional form is the simplest monotonically decreasing,
positive function with the correct limits. It parallels the
gravitational potential \(GM/r\), providing the ``attraction'' half of
the model.

\subsection{The Dissolution Cost
Function}\label{the-dissolution-cost-function}

Tight coupling also imposes \textbf{autonomy costs}: the HEC requires
conformity, extracts resources from the agent, constrains the agent's
action space (Definition 3 \citep{TC-I}), and exerts homogenizing
pressure on the agent's internal state. At sufficiently close coupling,
these costs overwhelm the benefits and erode the agent's boundary
integrity, leading to dissolution of identity.

We model the \textbf{dissolution cost rate} as:

\[C_{\text{dissolve}}(r) = \frac{\tau \, \mathcal{M}}{r^2}\]

where \(\tau > 0\) is the \textbf{dissolution coupling coefficient}, a
proportionality constant capturing the HEC's assimilation pressure per
unit mass.

\textbf{Properties:} - \(C_{\text{dissolve}}(r) > 0\) for all \(r > 0\).
- \(C_{\text{dissolve}}(r) \to \infty\) as \(r \to 0\), and
\textbf{faster} than \(E_{\text{in}}(r)\) (the \(1/r^2\) term dominates
\(1/r\) at small \(r\)). This is the critical structural feature: at
very close coupling, the costs always overwhelm the benefits. -
\(C_{\text{dissolve}}(r) \to 0\) as \(r \to \infty\) (distant agents are
not subject to assimilation pressure).

The \(1/r^2\) form ensures that the dissolution cost is a
\textbf{short-range dominant} force, negligible at large \(r\) but
overwhelming at small \(r\). This parallels the Lennard-Jones repulsion
in molecular physics and the ``tidal forces'' of the gravitational
analogy.

\textbf{Physical motivation:} The dissolution cost scales as \(1/r^2\)
rather than \(1/r\) because assimilation pressure involves not only
resource extraction but also \textbf{constraint imposition}: the HEC's
demand for conformity grows nonlinearly with coupling tightness. An
employee who works 60 hours/week for a corporation loses far more than
1.5× the autonomy of one working 40 hours. The constraints compound.

\subsection{The Entropy Maintenance
Baseline}\label{the-entropy-maintenance-baseline}

Independent of the HEC, the agent must continuously expend energy to
maintain its boundary against environmental entropy (Definition 1
\citep{TC-IV}):

\[C_{\text{maintain}} = \gamma_i \, B_i\]

where \(\gamma_i > 0\) is the entropy leakage rate and \(B_i > 0\) is
the boundary integrity. This is a \textbf{constant} cost: it does not
depend on \(r\) because entropy attacks the boundary regardless of the
agent's institutional affiliations. For the single-agent analysis that
follows, we write \(\gamma \equiv \gamma_i\) and suppress the index to
reduce notation.

\subsection{The Net Energy Rate}\label{the-net-energy-rate}

The agent's \textbf{net energy rate} at coupling distance \(r\) is:

\[\Pi(r) = E_{\text{in}}(r) - C_{\text{dissolve}}(r) - C_{\text{maintain}} = \frac{G \mathcal{M}}{r} - \frac{\tau \mathcal{M}}{r^2} - \gamma B_i\]

The agent survives if and only if \(\Pi(r) > 0\). The agent's goal is to
choose \(r\) to maximize \(\Pi(r)\), equivalently, to minimize the
\textbf{coexistence potential} \(V(r) = -\Pi(r)\).

\section{The Coexistence Potential}\label{the-coexistence-potential}

\subsection{Definition}\label{definition}

\begin{definition}[Coexistence Potential]
\label{def-coexistence-potential}

The \emph{coexistence potential} of an agent $i$ at coupling distance $r$ from a HEC with value mass $\mathcal{M}$ is:

$$V(r) = \frac{\tau \mathcal{M}}{r^2} - \frac{G \mathcal{M}}{r} + \gamma B_i$$

where $G > 0$ is the resource coupling coefficient, $\tau > 0$ is the dissolution coupling coefficient, $\gamma > 0$ is the entropy leakage rate, and $B_i > 0$ is the agent's boundary integrity. The potential is defined on $(0, \infty)$.
\end{definition}

The coexistence potential is the \textbf{negative} of the net energy
rate: \(V(r) = -\Pi(r)\). Regions where \(V(r) < 0\) are viable
(positive net energy); regions where \(V(r) > 0\) are non-viable (energy
deficit).

\subsection{Properties}\label{properties}

\begin{proposition}[Properties of the Coexistence Potential]
\label{prop-coexistence-potential-properties}

The coexistence potential $V(r)$ satisfies:

(a) $\lim_{r \to 0^+} V(r) = +\infty$ (dissolution dominance at close range).

(b) $\lim_{r \to \infty} V(r) = \gamma B_i > 0$ (starvation at large distance).

(c) $V$ is $C^\infty$ on $(0, \infty)$.

(d) $V$ has exactly one critical point on $(0, \infty)$, which is a global minimum.

(e) The minimum value $V_{\min}$ satisfies $V_{\min} < \gamma B_i$.
\end{proposition}

\begin{proof}
We prove each claim in turn:

(a) As $r \to 0^+$: $\tau \mathcal{M}/r^2 \to +\infty$, while $G\mathcal{M}/r \to +\infty$ at a slower rate. Since $1/r^2$ dominates $1/r$ as $r \to 0$, $V(r) \to +\infty$.

(b) As $r \to \infty$: both $\tau\mathcal{M}/r^2 \to 0$ and $G\mathcal{M}/r \to 0$, so $V(r) \to \gamma B_i > 0$.

(c) $V$ is a sum of rational functions of $r$ (with $r > 0$) and a constant, therefore smooth.

(d) Compute the first derivative:

$$V'(r) = -\frac{2\tau\mathcal{M}}{r^3} + \frac{G\mathcal{M}}{r^2} = \frac{\mathcal{M}}{r^3}\bigl(Gr - 2\tau\bigr)$$

Since $\mathcal{M} > 0$ and $r^3 > 0$, the sign of $V'(r)$ is determined by $(Gr - 2\tau)$:

\begin{itemize}
\item $V'(r) < 0$ for $r < 2\tau/G$ (potential decreasing)
\item $V'(r) = 0$ at $r = 2\tau/G$ (unique critical point)
\item $V'(r) > 0$ for $r > 2\tau/G$ (potential increasing)
\end{itemize}
The second derivative at the critical point:

$$V''(r) = \frac{6\tau\mathcal{M}}{r^4} - \frac{2G\mathcal{M}}{r^3}$$

At $r^* = 2\tau/G$:

$$V''(r^*) = \frac{6\tau\mathcal{M} \cdot G^4}{16\tau^4} - \frac{2G\mathcal{M} \cdot G^3}{8\tau^3} = \frac{G^4\mathcal{M}}{8\tau^3}\bigl(3 - 2\bigr) = \frac{G^4\mathcal{M}}{8\tau^3} > 0$$

Since $V''(r^*) > 0$, the critical point is a strict local minimum. By uniqueness of the critical point and the boundary behavior (parts a, b), it is the global minimum.

(e) $V(r^*) = \frac{\tau\mathcal{M}G^2}{4\tau^2} - \frac{G\mathcal{M}G}{2\tau} + \gamma B_i = \gamma B_i - \frac{G^2\mathcal{M}}{4\tau}$. Since $G^2\mathcal{M}/(4\tau) > 0$, we have $V(r^*) < \gamma B_i$.
\end{proof}

\subsection{The Coexistence Attractor}\label{the-coexistence-attractor}

\begin{definition}[Coexistence Attractor]
\label{def-coexistence-attractor}

The \emph{coexistence attractor} is the coupling distance that minimizes the coexistence potential:

$$r^* = \frac{2\tau}{G}$$

At this distance, the agent achieves the maximum net energy rate $\Pi(r^*) = G^2\mathcal{M}/(4\tau) - \gamma B_i$.
\end{definition}

\textbf{Interpretation:} The coexistence attractor \(r^*\) represents
the \textbf{optimal structural coupling} between the agent and the HEC.
It depends only on the ratio \(\tau/G\), the balance between dissolution
pressure and resource access, not on the HEC's mass \(\mathcal{M}\) or
the agent's boundary cost \(\gamma B_i\). This is a structural property:
the optimal degree of integration is determined by the \emph{nature} of
the coupling, not the \emph{scale} of the parties.

What does depend on \(\mathcal{M}\) is whether the agent can survive at
all (the depth of the potential well) and how wide the viable region is
(the Coexistence Band).

\subsection{The Potential Well Depth}\label{the-potential-well-depth}

\begin{definition}[Potential Well Depth]
\label{def-potential-well-depth}

The \emph{well depth} of the coexistence potential is:

$$D = \gamma B_i - V(r^*) = \frac{G^2 \mathcal{M}}{4\tau}$$

This is the maximum net energy surplus achievable by the agent when optimally coupled to the HEC.
\end{definition}

The well depth is proportional to \(\mathcal{M}\): a more massive center
creates a deeper well, supporting agents with higher maintenance costs.

\section{Gradient Dynamics and
Stability}\label{gradient-dynamics-and-stability}

\subsection{The Coupling Adjustment
Equation}\label{the-coupling-adjustment-equation}

Agents adjust their coupling distance toward higher net energy (lower
potential). The simplest model is \textbf{gradient dynamics}:

\[\dot{r} = -\mu \, V'(r) = -\frac{\mu \mathcal{M}}{r^3}(Gr - 2\tau)\]

where \(\mu > 0\) is the \textbf{adjustment rate} (mobility,
adaptability, i.e., how quickly the agent can change its institutional
affiliations).

\textbf{Dynamical behavior:}

\begin{itemize}
\tightlist
\item
  For \(r < r^*\): \(Gr - 2\tau < 0\), so \(\dot{r} > 0\): the agent
  moves \textbf{away} from the center (escaping dissolution pressure).
\item
  For \(r > r^*\): \(Gr - 2\tau > 0\), so \(\dot{r} < 0\): the agent
  moves \textbf{toward} the center (seeking resources).
\item
  At \(r = r^*\): \(\dot{r} = 0\): the agent is at equilibrium.
\end{itemize}

\subsection{Stability of the Coexistence
Attractor}\label{stability-of-the-coexistence-attractor}

\begin{theorem}[Stability of the Coexistence Attractor]
\label{thm-stability-coexistence-attractor}

Under the gradient dynamics $\dot{r} = -\mu V'(r)$, the coexistence attractor $r^* = 2\tau/G$ is:

(a) A fixed point: $\dot{r}|_{r=r^*} = 0$.

(b) Locally asymptotically stable: the eigenvalue $\lambda = -\mu V''(r^*) = -\mu G^4 \mathcal{M} / (8\tau^3) < 0$.

(c) Globally attracting on $(0, \infty)$: every trajectory starting at $r_0 \in (0, \infty)$ converges to $r^*$ as $t \to \infty$.
\end{theorem}

\begin{proof}
We prove each claim in turn:

(a) At $r^*$: $V'(r^*) = 0$, so $\dot{r} = 0$. ✓

(b) Linearize around $r^*$: $\dot{r} \approx -\mu V''(r^*)(r - r^*)$. From Proposition~\ref{prop-coexistence-potential-properties}, $V''(r^*) = G^4\mathcal{M}/(8\tau^3) > 0$. The eigenvalue of the linearized system is $\lambda = -\mu V''(r^*) < 0$. Hence $r^*$ is locally asymptotically stable.

(c) Define the Lyapunov function $\mathcal{W}(r) = V(r) - V(r^*)$. Then $\mathcal{W}(r) \geq 0$ for all $r > 0$ with equality iff $r = r^*$ (since $r^*$ is the unique global minimum of $V$). The time derivative along trajectories:

$$\dot{\mathcal{W}} = V'(r) \cdot \dot{r} = V'(r) \cdot \bigl(-\mu V'(r)\bigr) = -\mu \bigl(V'(r)\bigr)^2 \leq 0$$

with equality iff $V'(r) = 0$, i.e., $r = r^*$. By LaSalle's invariance principle, every trajectory converges to the largest invariant set within $\{r : \dot{\mathcal{W}} = 0\} = \{r^*\}$. Hence $r^*$ is globally attracting.

\textbf{Convergence rate:} Near $r^*$, the convergence is exponential with rate $|\lambda| = \mu G^4\mathcal{M}/(8\tau^3)$. Larger $\mu$ (more mobile agents) and larger $\mathcal{M}$ (more massive centers) accelerate convergence.
\end{proof}

\begin{corollary}[Resilience of the Equilibrium]
\label{cor-resilience-equilibrium}

The coexistence attractor is robust to perturbations: any temporary displacement from $r^*$ (due to external shocks, policy changes, economic disruptions) is self-correcting. The agent naturally returns to the optimal coupling distance.
\end{corollary}

\textbf{Physical interpretation:} A worker who is pushed closer to the
corporation than optimal (overwork, loss of autonomy) will naturally
seek more distance: reduced hours, unionization, job mobility. A worker
pushed too far (laid off, isolated) will naturally seek re-engagement:
job searching, networking. The gradient dynamics formalize this
behavioral regularity.

\section{The Stable Coexistence Band}\label{the-stable-coexistence-band}

\subsection{The Viability Condition}\label{the-viability-condition}

The agent survives if and only if \(\Pi(r) > 0\), equivalently
\(V(r) < 0\). The boundary of the viable region is defined by
\(V(r) = 0\):

\[\frac{\tau \mathcal{M}}{r^2} - \frac{G \mathcal{M}}{r} + \gamma B_i = 0\]

Multiplying by \(r^2\):

\[\gamma B_i \, r^2 - G\mathcal{M} \, r + \tau\mathcal{M} = 0\]

This is a quadratic in \(r\) with discriminant:

\[\Delta = G^2 \mathcal{M}^2 - 4\gamma B_i \tau \mathcal{M} = \mathcal{M}\bigl(G^2 \mathcal{M} - 4\gamma B_i \tau\bigr)\]

\subsection{The Minimum Viable HEC
Mass}\label{the-minimum-viable-hec-mass}

\begin{theorem}[Existence of the Coexistence Band]
\label{thm-existence-coexistence-band}

The Stable Coexistence Band exists if and only if the HEC's value mass exceeds a critical threshold:

$$\mathcal{M} > \mathcal{M}_{\min} \;\triangleq\; \frac{4\gamma B_i \tau}{G^2}$$

When $\mathcal{M} = \mathcal{M}_{\min}$, the open band is empty: $V(r) \geq 0$ everywhere, with equality only at $r^*$ — marginal, and non-surviving, since survival requires $\Pi(r) > 0$ strictly. When $\mathcal{M} < \mathcal{M}_{\min}$, no viable coupling distance exists: the agent cannot sustain itself near this HEC.
\end{theorem}

\begin{proof}
Two positive real roots of $V(r) = 0$ require: (i) $\Delta > 0$, which gives $\mathcal{M} > 4\gamma B_i \tau / G^2$; (ii) both roots positive, which is guaranteed since the sum $r_- + r_+ = G\mathcal{M}/(\gamma B_i) > 0$ and the product $r_- \cdot r_+ = \tau\mathcal{M}/(\gamma B_i) > 0$. When $\Delta = 0$, there is a single repeated root $r^* = G\mathcal{M}/(2\gamma B_i)$, and the minimum value $V(r^*) = 0$: the zero set of $V$ collapses to $\{r^*\}$ while the open band $\{V < 0\}$ is empty (marginal, non-surviving). When $\Delta < 0$, $V(r) > 0$ for all $r > 0$: the well is too shallow to reach negative values.
\end{proof}

\textbf{Physical interpretation:} A small corner shop (\(\mathcal{M}\)
low) cannot sustain employees who have high living costs (\(\gamma B_i\)
high). A massive technology corporation (\(\mathcal{M}\) high) can
sustain many agents with high boundary costs. The threshold
\(\mathcal{M}_{\min}\) is the minimum accumulated value at which a
center becomes a viable attractor.

\subsection{Inner and Outer
Boundaries}\label{inner-and-outer-boundaries}

\begin{definition}[Energetic Inner Boundary]
\label{def-dissolution-threshold}

The \emph{energetic inner boundary} is the inner boundary of the Stable Coexistence Band:

$$r_- = \frac{G\mathcal{M} - \sqrt{\Delta}}{2\gamma B_i}$$

Below $r_-$, the dissolution cost exceeds the resource benefit net of maintenance: $V(r) > 0$ for $r < r_-$. We reserve the term \emph{dissolution threshold} for the distinct boundary-degradation threshold $r_d$ of §\ref{boundary-dynamics-and-irreversibility}: $r_-$ marks where the energy balance turns negative, $r_d$ where active boundary erosion outruns repair.
\end{definition}

\begin{definition}[Starvation Threshold: Outer Boundary]
\label{def-starvation-threshold}

The \emph{starvation threshold} is the outer boundary of the Stable Coexistence Band:

$$r_+ = \frac{G\mathcal{M} + \sqrt{\Delta}}{2\gamma B_i}$$

Beyond $r_+$, the resource inflow is insufficient to cover maintenance and dissolution costs: $V(r) > 0$ for $r > r_+$.
\end{definition}

\begin{definition}[Stable Coexistence Band]
\label{def-stable-coexistence-band}

The \emph{Stable Coexistence Band} is the open interval:

$$\mathcal{B} = (r_-, r_+) = \left\{r > 0 : V(r) < 0\right\}$$

Within this band, the agent has positive net energy ($\Pi(r) > 0$) and can sustain its identity. The coexistence attractor $r^* = 2\tau / G$ lies within $\mathcal{B}$ whenever the band exists.
\end{definition}

\begin{lemma}[Attractor Containment]
\label{lem-attractor-containment}

If the Coexistence Band exists ($\mathcal{M} > \mathcal{M}_{\min}$), then $r_- < r^* < r_+$.
\end{lemma}

\begin{proof}
By Proposition~\ref{prop-coexistence-potential-properties}, $r^*$ is the unique global minimum of $V(r)$. When $\mathcal{M} > \mathcal{M}_{\min}$, $V(r^*) < 0$ while $V(r_-) = V(r_+) = 0$. Since $V$ is strictly decreasing on $(0, r^*)$ and strictly increasing on $(r^*, \infty)$, the zero $r_-$ must satisfy $r_- < r^*$ and the zero $r_+$ must satisfy $r_+ > r^*$.
\end{proof}

\subsection{Band Width: The Freedom
Bandwidth}\label{band-width-the-freedom-bandwidth}

\begin{theorem}[The Freedom Bandwidth Theorem]
\label{thm-freedom-bandwidth}

The width of the Stable Coexistence Band is:

$$w \;\triangleq\; r_+ - r_- = \frac{\sqrt{\Delta}}{\gamma B_i} = \frac{\sqrt{G^2 \mathcal{M}^2 - 4\gamma B_i \tau \mathcal{M}}}{\gamma B_i}$$

The bandwidth satisfies:

(a) $w > 0$ if and only if $\mathcal{M} > \mathcal{M}_{\min}$.

(b) $w$ is strictly increasing in $\mathcal{M}$: $\partial w / \partial \mathcal{M} > 0$.

(c) $w$ is strictly decreasing in each of $\gamma$, $B_i$, and $\tau$.

(d) $w$ is strictly increasing in $G$.

(e) In the limit $\mathcal{M} \gg \mathcal{M}_{\min}$: $w \approx G\mathcal{M}/(\gamma B_i)$, i.e., the bandwidth grows linearly with HEC mass.

(f) Near the viability threshold ($\mathcal{M} \to \mathcal{M}_{\min}^+$): $w \sim \sqrt{\mathcal{M} - \mathcal{M}_{\min}}$, i.e., the bandwidth vanishes as a square root.
\end{theorem}

\begin{proof}
From the quadratic formula: $r_\pm = (G\mathcal{M} \pm \sqrt{\Delta})/(2\gamma B_i)$, so $w = \sqrt{\Delta}/(\gamma B_i)$.

(a) $w > 0 \iff \Delta > 0 \iff \mathcal{M} > \mathcal{M}_{\min}$. ✓

(b) $\Delta = \mathcal{M}(G^2\mathcal{M} - 4\gamma B_i \tau) = G^2\mathcal{M}^2 - 4\gamma B_i \tau \mathcal{M}$. Then:

$$\frac{\partial \Delta}{\partial \mathcal{M}} = 2G^2\mathcal{M} - 4\gamma B_i \tau = 2\bigl(G^2\mathcal{M} - 2\gamma B_i \tau\bigr)$$

For $\mathcal{M} > \mathcal{M}_{\min} = 4\gamma B_i \tau/G^2$: $G^2\mathcal{M} > 4\gamma B_i \tau > 2\gamma B_i \tau$, so $\partial\Delta/\partial\mathcal{M} > 0$, therefore $\partial w/\partial\mathcal{M} > 0$. ✓

(c) Direct computation: $\partial\Delta/\partial\gamma = -4B_i\tau\mathcal{M} < 0$; $\partial\Delta/\partial B_i = -4\gamma\tau\mathcal{M} < 0$; $\partial\Delta/\partial\tau = -4\gamma B_i \mathcal{M} < 0$. Additionally, $w = \sqrt{\Delta}/(\gamma B_i)$, and $\gamma$, $B_i$ appear in the denominator, further decreasing $w$. ✓

(d) $\partial\Delta/\partial G = 2G\mathcal{M}^2 > 0$, so $\partial w / \partial G > 0$. ✓

(e) For $\mathcal{M} \gg \mathcal{M}_{\min}$: $\Delta \approx G^2\mathcal{M}^2$, so $w \approx G\mathcal{M}/(\gamma B_i)$. ✓

(f) Write $\Delta = G^2\mathcal{M}_{\min}(\mathcal{M} - \mathcal{M}_{\min}) + G^2(\mathcal{M} - \mathcal{M}_{\min})^2$. Near $\mathcal{M}_{\min}$, the linear term dominates: $\Delta \approx G^2\mathcal{M}_{\min}(\mathcal{M} - \mathcal{M}_{\min})$, so $w \approx G\sqrt{\mathcal{M}_{\min}(\mathcal{M} - \mathcal{M}_{\min})}/(\gamma B_i)$, confirming the square-root behavior.
\end{proof}

\begin{corollary}[Freedom Is Finite]
\label{cor-freedom-is-finite}

For any finite $\mathcal{M}$, the bandwidth $w$ is finite. No agent enjoys infinite freedom relative to any finite-mass center. Infinite freedom ($w \to \infty$) requires $\mathcal{M} \to \infty$ (an infinitely massive center) or $\gamma B_i \to 0$ (zero maintenance cost, i.e., an entity that requires no energy to exist, which is physically impossible).
\end{corollary}

\textbf{Physical interpretation of the Freedom Bandwidth:}

{\def\LTcaptype{none} % do not increment counter
\begin{longtable}[]{@{}
  >{\raggedright\arraybackslash}p{(\linewidth - 4\tabcolsep) * \real{0.3333}}
  >{\raggedright\arraybackslash}p{(\linewidth - 4\tabcolsep) * \real{0.3333}}
  >{\raggedright\arraybackslash}p{(\linewidth - 4\tabcolsep) * \real{0.3333}}@{}}
\toprule\noalign{}
\begin{minipage}[b]{\linewidth}\raggedright
Parameter
\end{minipage} & \begin{minipage}[b]{\linewidth}\raggedright
Effect on \(w\)
\end{minipage} & \begin{minipage}[b]{\linewidth}\raggedright
Interpretation
\end{minipage} \\
\midrule\noalign{}
\endhead
\bottomrule\noalign{}
\endlastfoot
\(\mathcal{M} \uparrow\) & \(w \uparrow\) & Richer centers offer wider
freedom (more room to maneuver) \\
\(G \uparrow\) & \(w \uparrow\) & More efficient resource coupling
expands viable distance \\
\(\gamma \uparrow\) & \(w \downarrow\) & Higher entropy pressure
(harsher environment) narrows freedom \\
\(B_i \uparrow\) & \(w \downarrow\) & More complex entities (higher
maintenance) have narrower bands \\
\(\tau \uparrow\) & \(w \downarrow\) & Stronger dissolution pressure
narrows freedom \\
\end{longtable}
}

\section{Boundary Dynamics and
Irreversibility}\label{boundary-dynamics-and-irreversibility}

The gradient dynamics established in the agent-center model
(§\ref{the-agent-center-model}) assume the agent's boundary integrity
\(B_i\) is constant. In reality, \(B_i\) is itself dynamic: it degrades
under assault from the HEC's assimilation pressure or from resource
deprivation, and it regenerates when the agent has surplus energy. The
boundary is not a passive reservoir that surplus refills at will. It is
an \textbf{active dissipative structure}: the machinery that performs
repair --- and that executes the coupling adjustments of the gradient
dynamics --- is itself part of the structure it maintains, so both
capacities degrade with the boundary. This section makes that structure
explicit and derives its central consequence: irreversibility is a
dynamical property of the coupled system, with an explicit basin of no
return, rather than an assumption.

\subsection{The Coupled Agent-Boundary
System}\label{the-coupled-agent-boundary-system}

\begin{definition}[Boundary Degradation Rate]
\label{def-boundary-degradation-rate}

The \emph{boundary degradation rate} due to HEC assimilation pressure at coupling distance $r$ is:

$$D_{\text{assimilate}}(r) = \frac{\sigma \mathcal{M}}{r^3}$$

where $\sigma > 0$ is the \textbf{assimilation intensity} coefficient. This is a short-range dominant force ($1/r^3$ decays faster than the dissolution cost $1/r^2$), reflecting that active boundary destruction requires even closer coupling than the autonomy costs defined in the dissolution cost function (§\ref{the-dissolution-cost-function}).
\end{definition}

\begin{definition}[Structural Capacity]
\label{def-structural-capacity}

Let $\bar{B} > 0$ denote the \textbf{reference integrity}: the embodied energy (Definition 6 \citep{TC-IV}) of the agent's fully intact boundary. Let $B_{\min} \in (0, \bar{B})$ denote the \textbf{collapse floor}: below $B_{\min}$, the boundary loses structural coherence, the machinery that performs repair and coupling adjustment is destroyed, and identity fails. The set $\{B_i \leq B_{\min}\}$ is \textbf{absorbing}: an agent whose integrity reaches the collapse floor is dissolved and ceases to exist. The \textbf{structural-capacity factor}

$$\zeta(B_i) = \max\left(\frac{B_i - B_{\min}}{\bar{B} - B_{\min}},\, 0\right) \in [0, 1]$$

is the fraction of full repair and mobility capacity the agent retains at integrity $B_i$: $\zeta$ is increasing on $[B_{\min}, \bar{B}]$ with $\zeta(B_{\min}) = 0$ and $\zeta(\bar{B}) = 1$. Throughout, $B_i$ ranges over $(B_{\min}, \bar{B}]$: repair restores the boundary toward its intact configuration and cannot raise $B_i$ above $\bar{B}$, so the dynamics saturate at the reference integrity.
\end{definition}

\begin{definition}[Boundary Regeneration Rate]
\label{def-boundary-regeneration-rate}

The \emph{boundary regeneration rate} is the integrity restored per unit time by the surplus energy the agent allocates to boundary repair:

$$R_{\text{repair}}(r, B_i) = \frac{\rho}{\kappa}\, \zeta(B_i)\, \max\bigl(\Pi(r), 0\bigr)$$

where $\rho \in (0, 1)$ is the \textbf{repair allocation fraction}, $\kappa = 1 + \kappa_{\text{clear}} + \kappa_{\text{pen}} + \kappa_{\text{irr}} > 1$ is the repair multiplier of \citep{TC-IV} (Definition 7 \citep{TC-IV}), and $\zeta(B_i)$ is the structural-capacity factor (Definition~\ref{def-structural-capacity}). Two mechanisms distinguish this from passive refilling of a reservoir. First, restoring $\Delta B$ of boundary costs $\kappa \cdot \Delta B$ of energy — the embodied energy must be re-supplied together with the Second-Law surcharges of rebuilding on a damaged site — so the allocated power $\rho \max(\Pi, 0)$ yields integrity at rate $(\rho/\kappa) \max(\Pi, 0)$. Second, repair is \textbf{autocatalytic}: it is performed by the boundary's own machinery, so repair capacity scales with the integrity that remains. A depleted structure cannot rebuild itself even when surplus energy is available: $R_{\text{repair}} \to 0$ as $B_i \to B_{\min}$ regardless of $\Pi$.
\end{definition}

Mobility is health-dependent for the same reason: adjusting the coupling
distance is work performed by the same structural machinery, so the
adjustment rate degrades with integrity,

\[\mu(B_i) = \mu_0 \, \zeta(B_i)\]

where \(\mu_0 > 0\) is the \textbf{baseline mobility}. (A stronger
variant would additionally have movement drain integrity, e.g.~a term
\(-c_{\text{move}}\,|\dot{r}|\) in \(\dot{B}_i\); the capacity factor
already captures the essential effect --- a depleted agent cannot afford
the migration --- and we adopt the minimal form.) The coupled dynamical
system governing coupling distance and boundary integrity is, for
\(B_i > B_{\min}\):

\[\dot{r} = -\mu_0\, \zeta(B_i)\, V'(r)\]

\[\dot{B}_i = \frac{\rho}{\kappa}\, \zeta(B_i)\, \max\bigl(\Pi(r), 0\bigr) - \frac{\sigma \mathcal{M}}{r^3} - \gamma B_i\]

with the absorbing dissolved state at \(B_i \leq B_{\min}\). Since
\(B_i\) is now a state variable, the net energy rate
\(\Pi(r) = G\mathcal{M}/r - \tau\mathcal{M}/r^2 - \gamma B_i\) and every
quantity defined from it --- \(V\), the band boundaries \(r_\pm\),
\(\mathcal{M}_{\min}\) --- are read at the current integrity; the
attractor \(r^* = 2\tau/G\) is independent of \(B_i\). The
constant-mobility gradient dynamics of the stability analysis
(§\ref{gradient-dynamics-and-stability}) are recovered in the healthy
regime: for an agent whose integrity holds at the equilibrium \(B_i^*\)
derived below, the first equation is the earlier flow with
\(\mu = \mu_0 \zeta(B_i^*)\), a constant time-rescaling that leaves
every phase-portrait conclusion (attractor location, stability, band
structure) unchanged.

\subsection{The Healthy Equilibrium and the Basin of No
Return}\label{the-healthy-equilibrium-and-the-basin-of-no-return}

The coupled system must pass two checks. It must sustain viable agents
--- a healthy agent in the band should hold a stable integrity strictly
above the collapse floor --- and it should locate exactly where
viability fails. Both follow from the fixed-\(r\) boundary dynamics,
which are quadratic in \(B_i\).

\begin{proposition}[Healthy Boundary Equilibrium]
\label{prop-healthy-boundary-equilibrium}

Fix $r$ and define the floor-level net rate $\Pi_0(r) \triangleq G\mathcal{M}/r - \tau\mathcal{M}/r^2 - \gamma B_{\min}$, the floor-level drain $d_0(r) \triangleq \sigma\mathcal{M}/r^3 + \gamma B_{\min}$, and the repair-capacity density $q \triangleq \rho / \bigl(\kappa(\bar{B} - B_{\min})\bigr)$. Suppose the \textbf{repair-viability condition} holds at $r$:

$$q\, \Pi_0(r) \;>\; \gamma + 2\sqrt{q \gamma\, d_0(r)}.$$

Then $\dot{B}_i = 0$ at fixed $r$ has exactly two roots $B_{\mathrm{u}}(r) < B_i^*(r)$ in $(B_{\min}, \infty)$, given by $B_{\min} + x_\mp$ with

$$x_\mp = \frac{\bigl(q \Pi_0(r) - \gamma\bigr) \mp \sqrt{\bigl(q \Pi_0(r) - \gamma\bigr)^2 - 4 q \gamma\, d_0(r)}}{2 q \gamma},$$

and:

(a) $B_i^*(r) = B_{\min} + x_+$ is locally asymptotically stable for the boundary dynamics at fixed $r$, with $\Pi(r) > 0$ at $B_i = B_i^*(r)$. If $x_+ > \bar{B} - B_{\min}$, the equilibrium saturates at the reference integrity $\bar{B}$.

(b) $B_{\mathrm{u}}(r)$ is unstable: $\dot{B}_i < 0$ for $B_i \in (B_{\min}, B_{\mathrm{u}}(r))$ and $\dot{B}_i > 0$ for $B_i \in \bigl(B_{\mathrm{u}}(r), \min(B_i^*(r), \bar{B})\bigr)$.

(c) If the repair-viability condition holds at $r = r^*$ with $B_{\mathrm{u}}(r^*) < \bar{B}$ — the \textbf{viable-agent condition} — then $\bigl(r^*, \min(B_i^*(r^*), \bar{B})\bigr)$ is a locally asymptotically stable equilibrium of the coupled system: a healthy agent at the attractor maintains a boundary equilibrium strictly above the collapse floor, and the attractor and band results of the preceding sections apply verbatim with $B_i = B_i^*$ and $\mu = \mu_0 \zeta(B_i^*)$.
\end{proposition}

\begin{proof}
Write $x = B_i - B_{\min} > 0$, so $\zeta(B_i) = x / (\bar{B} - B_{\min})$ and $\Pi(r) = \Pi_0(r) - \gamma x$. On the branch $\Pi > 0$ (i.e. $x < \Pi_0(r)/\gamma$), the fixed-$r$ boundary dynamics read

$$\dot{B}_i = q\, x \bigl(\Pi_0(r) - \gamma x\bigr) - \gamma x - d_0(r) = -q\gamma x^2 + \bigl(q\Pi_0(r) - \gamma\bigr) x - d_0(r),$$

a downward-opening parabola in $x$ with value $-d_0(r) < 0$ at $x = 0$. It has two positive roots $x_- < x_+$ if and only if $q\Pi_0(r) - \gamma > 0$ and the discriminant $(q\Pi_0(r) - \gamma)^2 - 4q\gamma d_0(r)$ is positive — jointly, the repair-viability condition. The parabola is negative on $(0, x_-)$, positive on $(x_-, x_+)$, and negative beyond $x_+$, which gives the sign pattern in (b); its slope at $x_\pm$ is $\mp\sqrt{(q\Pi_0(r) - \gamma)^2 - 4q\gamma d_0(r)}$, so $B_i^*(r)$ is locally asymptotically stable and $B_{\mathrm{u}}(r)$ unstable. At $x = \Pi_0(r)/\gamma$ (where $\Pi = 0$) the parabola evaluates to $-\Pi_0(r) - d_0(r) < 0$; since this point lies to the right of the vertex $\bigl(q\Pi_0(r) - \gamma\bigr)/(2q\gamma) < \Pi_0(r)/(2\gamma)$, the negative value places it beyond the right root, hence $x_+ < \Pi_0(r)/\gamma$ and $\Pi > 0$ at $B_i^*(r)$; for $x \geq \Pi_0(r)/\gamma$ the repair term vanishes and $\dot{B}_i = -\sigma\mathcal{M}/r^3 - \gamma B_i < 0$, so there are no further fixed points. If $x_+ > \bar{B} - B_{\min}$, then $\dot{B}_i > 0$ on $(B_{\mathrm{u}}(r), \bar{B})$ and the saturated state $\bar{B}$ is attracting from below.

For (c): $V'(r^*) = 0$ independently of $B_i$, so $\dot{r} = 0$ on the line $r = r^*$ and $\partial \dot{r}/\partial B_i = -\mu_0\, \zeta'(B_i)\, V'(r^*) = 0$ there. The Jacobian at $(r^*, B_i^*)$ is therefore lower-triangular with eigenvalues $-\mu_0 \zeta(B_i^*) V''(r^*) < 0$ (Proposition~\ref{prop-coexistence-potential-properties}) and $-\sqrt{(q\Pi_0(r^*) - \gamma)^2 - 4q\gamma d_0(r^*)} < 0$; both negative, so the equilibrium is locally asymptotically stable (in the saturated case, $\bar{B}$ is attracting from below by the sign of $\dot{B}_i$). With $B_i$ held at $B_i^*$, the flow $\dot{r} = -\mu_0 \zeta(B_i^*) V'(r)$ is the gradient flow of Theorem~\ref{thm-stability-coexistence-attractor} under the constant rescaling $\mu = \mu_0 \zeta(B_i^*) > 0$, which preserves trajectories and their limits; the band, bandwidth, and comparative-statics results depend only on $V$ and hold with $B_i = B_i^*$.
\end{proof}

The unstable root \(B_{\mathrm{u}}(r)\) is the fixed-\(r\) boundary of
the recovery region. The next theorem extracts its
trajectory-independent core: a critical integrity below which no motion
in \(r\) --- and no future surplus --- can save the agent.

\begin{theorem}[The Basin of No Return]
\label{thm-basin-of-no-return}

Define the \textbf{surplus ceiling} $\bar{\Pi} \triangleq G^2\mathcal{M}/(4\tau) - \gamma B_{\min}$: the largest net energy rate available at any coupling distance to an agent at the collapse floor (assume $\bar{\Pi} > 0$, which holds whenever the band is nonempty for some admissible integrity). Let

$$B_c \;\triangleq\; B_{\min} + \frac{\bigl(q\bar{\Pi} - \gamma\bigr) - \sqrt{\bigl(q\bar{\Pi} - \gamma\bigr)^2 - 4 q \gamma^2 B_{\min}}}{2 q \gamma},$$

the smaller root of $q\,(B_i - B_{\min})\bigl(G^2\mathcal{M}/(4\tau) - \gamma B_i\bigr) = \gamma B_i$. If the repair-viability condition of Proposition~\ref{prop-healthy-boundary-equilibrium} holds at some $r$, then $B_c$ exists and $B_{\min} < B_c \leq B_{\mathrm{u}}(r)$; if moreover the agent admits a healthy equilibrium — the viable-agent condition of Proposition~\ref{prop-healthy-boundary-equilibrium}(c) — then $B_c \leq B_{\mathrm{u}}(r^*) < \bar{B}$, so the basin is a proper layer of the admissible range. Furthermore:

(a) **(Collapse below $B_c$.)** If $B_i(t_0) < B_c$, then along every trajectory of the coupled system $\dot{B}_i < 0$ for all $t \geq t_0$ *regardless of the path $r(t)$*, and $B_i$ reaches the collapse floor $B_{\min}$ in finite time. No later reversal is possible even if $\Pi > 0$ subsequently: the estimate in the proof already grants the agent the maximal surplus $G^2\mathcal{M}/(4\tau) - \gamma B_i$ at every instant.

(b) \textbf{(Mobility death.)} As $B_i \downarrow B_{\min}$, $\mu(B_i) = \mu_0 \zeta(B_i) \to 0$: the collapsing agent progressively loses the capacity to relocate, and the dissolved state is reached at asymptotically frozen coupling distance.

(c) \textbf{(Recovery above the basin.)} If $B_i(t_0) > B_{\mathrm{u}}(r^*)$ and $r(t_0) = r^*$, then $\dot{B}_i > 0$ and the agent converges to the healthy equilibrium of Proposition~\ref{prop-healthy-boundary-equilibrium}(c). More generally, $B_c$ is a certified inner bound of the basin: below $B_c$, collapse is unconditional; between $B_c$ and $B_{\mathrm{u}}(r(t))$, the outcome depends on the full trajectory.
\end{theorem}

\begin{proof}
For $B_i \in (B_{\min}, \bar{B}]$ we have $\zeta(B_i) \leq 1$ and, for every $r > 0$, $\max(\Pi(r), 0) \leq \max\bigl(G^2\mathcal{M}/(4\tau) - \gamma B_i,\, 0\bigr)$ — the maximum of $\Pi$ over $r$ at the current integrity, attained at $r^*$. Dropping the assimilation drain $-\sigma\mathcal{M}/r^3 \leq 0$:

$$\dot{B}_i \;\leq\; h(B_i) \;\triangleq\; q\,(B_i - B_{\min})\left(\frac{G^2\mathcal{M}}{4\tau} - \gamma B_i\right) - \gamma B_i,$$

valid wherever $G^2\mathcal{M}/(4\tau) - \gamma B_i \geq 0$, in particular on $(B_{\min}, B_c]$ (verified below). In $x = B_i - B_{\min}$, $h = -q\gamma x^2 + (q\bar{\Pi} - \gamma)x - \gamma B_{\min}$, a downward parabola with $h(B_{\min}) = -\gamma B_{\min} < 0$. Since $h$ dominates the fixed-$r$ dynamics of Proposition~\ref{prop-healthy-boundary-equilibrium} (it omits the drain $\sigma\mathcal{M}/r^3 \geq 0$), $h(B_{\mathrm{u}}(r)) \geq 0$ whenever repair-viability holds at $r$; hence $h$ has real roots and its smaller root $B_c$ satisfies $B_{\min} < B_c \leq B_{\mathrm{u}}(r)$. Under the viable-agent condition this holds at $r = r^*$, where additionally $B_{\mathrm{u}}(r^*) < \bar{B}$, giving $B_c < \bar{B}$. At $B_c$, $h(B_c) = 0$ with $q(B_c - B_{\min}) > 0$ forces $G^2\mathcal{M}/(4\tau) - \gamma B_c > 0$, so the bound applies on all of $(B_{\min}, B_c]$.

(a) On $(B_{\min}, B_c)$, $h < 0$ and $h$ is increasing (the parabola's vertex lies beyond $B_c$). Let $B_i(t_0) = B_0 < B_c$. Then $\dot{B}_i \leq h(B_i) < 0$, so $B_i$ decreases and remains below $B_0$; consequently $\dot{B}_i \leq h(B_i(t)) \leq h(B_0) < 0$ for all $t \geq t_0$, and $B_i$ reaches $B_{\min}$ no later than $t_0 + (B_0 - B_{\min})/|h(B_0)|$. The estimate used the global maximum of the surplus and ignored assimilation, so it holds along every path $r(t)$ and is unaffected by any later sign change of $\Pi$.

(b) Immediate from $\mu(B_i) = \mu_0 \zeta(B_i)$ and $\zeta(B_{\min}) = 0$.

(c) On the line $r = r^*$ the flow has $\dot{r} = 0$ for all $B_i$ (since $V'(r^*) = 0$), so the fixed-$r^*$ analysis of Proposition~\ref{prop-healthy-boundary-equilibrium} governs: $\dot{B}_i > 0$ on $\bigl(B_{\mathrm{u}}(r^*), \min(B_i^*(r^*), \bar{B})\bigr)$ and $B_i \to \min(B_i^*(r^*), \bar{B})$.
\end{proof}

\emph{Remark.} When \(q\bar{\Pi} \gg \gamma\), the smaller root
satisfies
\(B_c \approx B_{\min}\bigl(1 + \gamma/(q\bar{\Pi} - \gamma)\bigr)\):
for agents with a large surplus ceiling, the basin is a thin layer above
the collapse floor, and it thickens as \(\bar{\Pi}\) falls toward
\(\gamma/q = \gamma\kappa(\bar{B} - B_{\min})/\rho\). The hypothesis
\(B_i < B_c\) is produced by the two failure regimes analyzed next:
sustained coupling below the dissolution threshold, and sustained energy
deficit outside the band.

\subsection{The Dissolution
Catastrophe}\label{the-dissolution-catastrophe}

\begin{theorem}[Irreversibility of Dissolution]
\label{thm-irreversibility-dissolution}

Define the \textbf{dissolution threshold}

$$r_d \;\triangleq\; \left(\frac{\kappa\, \sigma \mathcal{M}}{\rho\, \bar{\Pi} + \kappa \gamma B_{\min}}\right)^{1/3}, \qquad \text{equivalently} \qquad \frac{\sigma\mathcal{M}}{r_d^3} = \frac{\rho}{\kappa}\bar{\Pi} + \gamma B_{\min},$$

where $\bar{\Pi} = G^2\mathcal{M}/(4\tau) - \gamma B_{\min}$ is the surplus ceiling of Theorem~\ref{thm-basin-of-no-return}. Then:

(a) \textbf{(Uniform decline.)} Whenever $r(t) < r_d$, every admissible integrity $B_i \in (B_{\min}, \bar{B}]$ satisfies $\dot{B}_i < -\gamma\bigl(B_i + B_{\min}\bigr) < 0$.

(b) \textbf{(Conditional finite-time dissolution.)} If $r(t) < r_d$ for all $t \geq t_0$ — in particular if the coupling distance is held below $r_d$, or if $r_d \geq r^*$ by (c) — then $B_i$ reaches the collapse floor at a finite time $t_d \leq t_0 + \bigl(B_i(t_0) - B_{\min}\bigr)/(2\gamma B_{\min})$: the agent dissolves.

(c) \textbf{(Irreversibility.)} If $r_d \geq r^*$, every trajectory with $r(t_0) < r_d$ satisfies $r(t) < r_d$ for all $t \geq t_0$ — the attractor itself lies inside the dissolution zone — and (b) applies unconditionally. If $r_d < r^*$ and the penetration depth exceeds the critical depth of the remark below ($\ell > \ell_c$), then $B_i$ falls below $B_c$ before the agent can escape $\{r < r_d\}$, and collapse is irreversible by Theorem~\ref{thm-basin-of-no-return} regardless of the subsequent trajectory.
\end{theorem}

\begin{proof}
(a) For $B_i \in (B_{\min}, \bar{B}]$, $\zeta(B_i) \leq 1$ and $\max(\Pi(r), 0) \leq \max\bigl(G^2\mathcal{M}/(4\tau) - \gamma B_i,\, 0\bigr) \leq \bar{\Pi}$ (using $\bar{\Pi} > 0$), so $R_{\text{repair}} \leq (\rho/\kappa)\bar{\Pi}$. For $r < r_d$, $D_{\text{assimilate}}(r) = \sigma\mathcal{M}/r^3 > \sigma\mathcal{M}/r_d^3 = (\rho/\kappa)\bar{\Pi} + \gamma B_{\min}$. Hence

$$\dot{B}_i < \frac{\rho}{\kappa}\bar{\Pi} - \left(\frac{\rho}{\kappa}\bar{\Pi} + \gamma B_{\min}\right) - \gamma B_i = -\gamma\bigl(B_i + B_{\min}\bigr).$$

(b) While $r(t) < r_d$ and $B_i > B_{\min}$, part (a) gives $\dot{B}_i < -2\gamma B_{\min}$, a decline bounded away from zero; integrating, $B_i$ falls from $B_i(t_0)$ to $B_{\min}$ within $\bigl(B_i(t_0) - B_{\min}\bigr)/(2\gamma B_{\min})$.

(c) Suppose $r_d \geq r^*$. The sign of $\dot{r} = -\mu_0\zeta(B_i)V'(r)$ is that of $-(Gr - 2\tau)$: positive for $r < r^*$, negative for $r > r^*$, zero at $r^*$. A trajectory with $r(t_0) \leq r^*$ increases toward $r^*$ and remains $\leq r^* \leq r_d$; one with $r(t_0) \in (r^*, r_d)$ decreases toward $r^*$ and remains in $(r^*, r(t_0)]$. Either way $r(t) < r_d$ for all $t \geq t_0$, and (b) applies. The case $r_d < r^*$ is the race condition resolved in the following remark.
\end{proof}

\emph{Remark (Race-Condition Closure).}
Theorem~\ref{thm-irreversibility-dissolution} establishes dissolution
conditional on remaining below \(r_d\). When \(r_d < r^*\), the gradient
dynamics push \(r\) upward toward \(r^*\) on \((0, r_d)\), so we must
verify that boundary depletion outpaces escape. We resolve this via the
phase-space derivative \(dB_i/dr = \dot{B}_i/\dot{r}\) along
trajectories of the coupled system
(§\ref{the-coupled-agent-boundary-system}) with \(r \in (0, r_d)\),
where \(2\tau - Gr > 0\) and
\(\dot{r} = \mu_0\zeta(B_i)\mathcal{M}(2\tau - Gr)/r^3 > 0\):

\[\frac{dB_i}{dr} = \frac{\dfrac{\rho}{\kappa}\zeta(B_i)\max\bigl(\Pi(r),0\bigr) - \dfrac{\sigma\mathcal{M}}{r^3} - \gamma B_i}{\mu_0\,\zeta(B_i)\,\mathcal{M}(2\tau - Gr)/r^3}.\]

Both the assimilation drain and the escape velocity share the
\(\mathcal{M}/r^3\) scaling, which cancels in the ratio. Bound the
numerator using \(\max(\Pi(r), 0) \leq \bar{\Pi}\) (valid at every
admissible integrity, since \(\bar{\Pi}\) is evaluated at the collapse
floor) and \(-\gamma B_i \leq 0\); multiplying by \(r^3/\mathcal{M}\)
and using
\(r^3 \leq r_d^3 = \kappa\sigma\mathcal{M}/(\rho\bar{\Pi} + \kappa\gamma B_{\min})\):

\[\frac{\rho}{\kappa}\,\zeta(B_i)\,\bar{\Pi}\,\frac{r^3}{\mathcal{M}} - \sigma \;\leq\; \sigma\left(\frac{\rho\bar{\Pi}}{\rho\bar{\Pi} + \kappa\gamma B_{\min}}\,\zeta(B_i) - 1\right) \;\leq\; -\,\frac{\sigma\kappa\gamma B_{\min}}{\rho\bar{\Pi} + \kappa\gamma B_{\min}},\]

where the last step uses \(\zeta(B_i) \leq 1\). Applying the
conservative envelopes \(\zeta(B_i) \leq 1\) and
\(2\tau - Gr \leq 2\tau\) in the denominator yields the uniform bound
\(dB_i/dr \leq -\alpha\) where

\[\alpha \;\triangleq\; \frac{\sigma\kappa\gamma B_{\min}}{2\mu_0\tau\,\bigl(\rho\bar{\Pi} + \kappa\gamma B_{\min}\bigr)} \;=\; \frac{\gamma B_{\min}\, r_d^3}{2\mu_0 \tau\, \mathcal{M}} \;>\; 0.\]

Integrating from \(r_0\) to \(r\) gives
\(B_i(r) \leq B_i^{(0)} - \alpha(r - r_0)\). Define the \textbf{critical
penetration depth}

\[\ell_c \;\triangleq\; \frac{B_i^{(0)} - B_c}{\alpha} = \frac{2\mu_0\tau\,\bigl(\rho\bar{\Pi} + \kappa\gamma B_{\min}\bigr)\bigl(B_i^{(0)} - B_c\bigr)}{\sigma\kappa\gamma B_{\min}}.\]

If the agent starts at \(r_0\) with penetration depth
\(\ell \triangleq r_d - r_0 > \ell_c\), then \(B_i\) falls below the
basin boundary \(B_c\) before \(r\) can reach \(r_d\); by
Theorem~\ref{thm-basin-of-no-return}, collapse to \(B_{\min}\) then
completes in finite time regardless of the subsequent trajectory --- the
race is lost. The bound is conservative on three counts: the envelope
\(2\tau - Gr \leq 2\tau\), the omission of the maintenance drain
\(\gamma B_i\), and above all the substitution \(\zeta(B_i) \leq 1\) in
the escape velocity. Along the true trajectory the decline per unit
escape distance grows without bound as \(B_i \to B_{\min}\), because
mobility dies with the boundary (\(\mu(B_i) \to 0\)) while assimilation
persists. The condition \(\ell > \ell_c\) is therefore sufficient but
far from necessary for irreversibility.

Physical interpretation: an employee drawn into a corporation's inner
circle eventually crosses a depth from which the escape gradient, always
present, cannot outpace identity erosion. The critical depth \(\ell_c\)
scales with baseline mobility \(\mu_0\) (mobile agents tolerate deeper
excursions) and inversely with assimilation intensity \(\sigma\)
(stronger assimilators set shallower traps).

\begin{corollary}[The Assimilation Trap]
\label{cor-assimilation-trap}

The dissolution threshold $r_d$ and the energetic inner boundary $r_-$ (Definition~\ref{def-dissolution-threshold}) are distinct thresholds, and either ordering can occur. If $r_d > r_-$, the dissolution zone extends into the viable band: an agent at $r \in (r_-, r_d)$ has positive net energy ($\Pi(r) > 0$) while its boundary degrades ($\dot{B}_i < 0$, by Theorem~\ref{thm-irreversibility-dissolution}a). Moreover, as $B_i$ falls the inner band edge $r_- = (G\mathcal{M} - \sqrt{\Delta})/(2\gamma B_i)$ moves inward (from $r^*$ toward $\tau/G$ as $\gamma B_i \to 0$), so an agent that entered the dissolution zone below $r_-$ becomes energetically viable ($r > r_-$) while remaining inside the dissolution zone ($r < r_d$). This is the \textbf{assimilation trap}: the agent appears viable by energy metrics while losing its identity — irreversibly so once $B_i < B_c$ (Theorem~\ref{thm-basin-of-no-return}).
\end{corollary}

\textbf{Physical interpretation:} An employee (agent) at a
mega-corporation (HEC) may earn a good salary (positive \(\Pi(r)\))
while simultaneously losing creative autonomy, decision-making capacity,
and identity (\(B_i\) declining). By the time the boundary has degraded
enough for the agent to ``feel'' the loss, the process may already be
irreversible: once \(B_i\) has fallen below the basin boundary \(B_c\),
the agent has been assimilated (corporate culture has overwritten
individual identity). Similarly, a small nation within a superpower's
sphere of influence may experience economic benefits while losing
sovereignty.

\subsection{The Starvation Spiral}\label{the-starvation-spiral}

\begin{theorem}[The Starvation Spiral]
\label{thm-starvation-spiral}

Suppose the agent is outside the Coexistence Band, $r(t) > r_+$, so that $\Pi(r) < 0$. Then:

(a) The agent runs an energy deficit and allocates nothing to repair: $R_{\text{repair}} = 0$.

(b) $\dot{B}_i = -D_{\text{assimilate}}(r) - \gamma B_i < 0$: integrity declines monotonically, with $B_i(t) \leq B_i^{(0)} e^{-\gamma t}$ for as long as the agent remains outside the band. In the isolation limit ($r \gg r_+$, assimilation negligible), the decline is exponential and reaches the collapse floor at the finite time

$$t^* = \frac{1}{\gamma} \ln \frac{B_i^{(0)}}{B_{\min}}$$

— the exponential never reaches zero, but dissolution requires only reaching $B_{\min} > 0$.

(c) \textbf{(Race between recovery and depletion.)} Two recovery channels operate while the agent starves: the inward drift $\dot{r} = -\mu_0\zeta(B_i)V'(r) < 0$ (for $r > r^*$) carries the agent back toward the band, and the falling maintenance load moves the outer boundary $r_+(B_i) = (G\mathcal{M} + \sqrt{\Delta})/(2\gamma B_i)$ outward toward the agent. While $B_i > B_c$ the expanding boundary remains inside $r_+(B_c)$, so re-entry from $r_0$ takes at least $\bigl(r_0^3 - r_+(B_c)^3\bigr)/(3\mu_0\mathcal{M}G)$, while $B_i$ falls below the basin boundary $B_c$ no later than $t_c = (1/\gamma)\ln\bigl(B_i^{(0)}/B_c\bigr)$ if the deficit persists. Both channels close at $B_c$: below it, Theorem~\ref{thm-basin-of-no-return} applies regardless of position or subsequent surplus.

(d) \textbf{(Irreversible starvation for depleted agents.)} If $B_i < B_c$ — in particular whenever the agent starts sufficiently far out or starves long enough; a sufficient condition is

$$r_0^3 \;\geq\; r_+(B_c)^3 + \frac{3\mu_0\mathcal{M}G}{\gamma}\ln\frac{B_i^{(0)}}{B_c},$$

where $r_+(B_c)$ is the outer boundary evaluated at integrity $B_c$ — then repair capacity and mobility vanish together ($\zeta \to 0$), the inward drift stalls before re-entry, and $B_i$ reaches $B_{\min}$ in finite time: starvation is irreversible. Conversely, a briefly starved agent whose integrity remains above the basin can recover by re-entering the band (Proposition~\ref{prop-healthy-boundary-equilibrium}): starvation becomes fatal not when the band is exited but when depletion crosses $B_c$.
\end{theorem}

\begin{proof}
(a) Outside $\mathcal{B}$: $V(r) > 0$, so $\Pi(r) < 0$ and $\max(\Pi(r), 0) = 0$.

(b) $\dot{B}_i = -\sigma\mathcal{M}/r^3 - \gamma B_i \leq -\gamma B_i$; Grönwall's inequality gives $B_i(t) \leq B_i^{(0)} e^{-\gamma t}$ while the deficit lasts. In the isolation limit $\sigma\mathcal{M}/r^3 \to 0$, the inequality is an equality and $B_i(t) = B_i^{(0)} e^{-\gamma t}$ reaches $B_{\min}$ exactly at $t^*$.

(c) For $r > r^*$, $0 < V'(r) = \mathcal{M}(Gr - 2\tau)/r^3 \leq \mathcal{M}G/r^2$ and $\zeta(B_i) \leq 1$, so $|\dot{r}| \leq \mu_0\mathcal{M}G/r^2$; for $r \leq r^*$, $\dot{r} \geq 0$. In either case $\frac{d}{dt}\,r^3 \geq -3\mu_0\mathcal{M}G$, hence $r(t)^3 \geq r_0^3 - 3\mu_0\mathcal{M}G\,t$: reaching any target $R < r_0$ takes at least $(r_0^3 - R^3)/(3\mu_0\mathcal{M}G)$. Since $r_+(B_i)$ is decreasing in $B_i$, the moving boundary satisfies $r_+(B_i(t)) < r_+(B_c)$ while $B_i(t) > B_c$, so re-entry before basin crossing requires reaching $R = r_+(B_c)$. The depletion bound follows from (b).

(d) If $B_i^{(0)} \leq B_c$, Theorem~\ref{thm-basin-of-no-return}(a) applies immediately, so assume $B_i^{(0)} > B_c$ and the displayed condition. Let $t_b$ be the first time the deficit fails ($\Pi = 0$) before $B_i$ reaches $B_c$, if any such time exists. For $t < t_b$, (b) gives $B_i(t) \leq B_i^{(0)} e^{-\gamma t}$ and (c) gives $r(t)^3 \geq r_0^3 - 3\mu_0\mathcal{M}G\,t$. At $t_b$, $\gamma B_i(t_b) = G\mathcal{M}/r(t_b) - \tau\mathcal{M}/r(t_b)^2$ with $B_i(t_b) > B_c$, which requires $r(t_b) < r_+(B_c)$ (the right-hand side equals $\gamma B_c$ at $r = r_+(B_c)$ and is below $\gamma B_c$ for all $r > r_+(B_c)$). Hence $r_0^3 - 3\mu_0\mathcal{M}G\,t_b < r_+(B_c)^3$, so the displayed condition forces $t_b > t_c = (1/\gamma)\ln\bigl(B_i^{(0)}/B_c\bigr)$ — but then $B_i(t_b) \leq B_i^{(0)} e^{-\gamma t_b} < B_c$, a contradiction. The deficit therefore persists at least until $B_i \leq B_c$, and Theorem~\ref{thm-basin-of-no-return}(a) completes the collapse to $B_{\min}$ in finite time, with mobility dying en route (Theorem~\ref{thm-basin-of-no-return}b).
\end{proof}

\textbf{Physical interpretation:} An entity that completely cuts itself
off from all energy centers (autarky, isolationism) depletes its stores;
without resource inflow, entropy wins, and the boundary reaches its
collapse floor in finite time. Transient starvation is survivable: an
agent that re-engages while its integrity remains above the basin
boundary \(B_c\) can rebuild. What is fatal is prolonged deficit --- by
the time integrity has fallen below \(B_c\), the capacities needed for
recovery (mobility and repair machinery) have themselves been consumed.
This maps to isolated communities, autarkic states, or any
self-maintaining system severed from its resource base past the point
where re-engagement is still within its remaining means.

\section{Comparative Statics and
Inequality}\label{comparative-statics-and-inequality}

\subsection{Effect of Value Mass on the Coexistence
Band}\label{effect-of-value-mass-on-the-coexistence-band}

\begin{proposition}[Value Mass Comparative Statics]
\label{prop-value-mass-comparative-statics}

As the HEC's value mass $\mathcal{M}$ increases:

(a) The well depth $D = G^2\mathcal{M}/(4\tau)$ increases linearly.

(b) The bandwidth $w$ increases (Theorem~\ref{thm-freedom-bandwidth}b).

(c) The energetic inner boundary $r_-$ decreases (moves inward):

$$\frac{\partial r_-}{\partial \mathcal{M}} = \frac{1}{2\gamma B_i}\left(G - \frac{G^2\mathcal{M} - 2\gamma B_i \tau}{\sqrt{\Delta}}\right) < 0 \quad \text{for } \mathcal{M} > \mathcal{M}_{\min}$$

(d) The starvation threshold $r_+$ increases (the outer boundary moves outward):

$$\frac{\partial r_+}{\partial \mathcal{M}} = \frac{1}{2\gamma B_i}\left(G + \frac{G^2\mathcal{M} - 2\gamma B_i \tau}{\sqrt{\Delta}}\right) > 0$$

(e) The coexistence attractor $r^* = 2\tau/G$ is unchanged.
\end{proposition}

\textbf{Physical interpretation:} More massive centers support wider
freedom. A wealthier, more institutionally robust state provides its
citizens with both more protection against starvation (larger \(r_+\))
and more residual autonomy before assimilation (smaller \(r_-\)). The
optimal coupling structure (\(r^*\)) is invariant; what changes is the
margin of error.

\subsection{Effect of Agent Boundary
Integrity}\label{effect-of-agent-boundary-integrity}

\begin{proposition}[Boundary Integrity Comparative Statics]
\label{prop-boundary-integrity-comparative-statics}

As the agent's boundary integrity $B_i$ increases:

(a) The minimum viable HEC mass $\mathcal{M}_{\min} = 4\gamma B_i \tau / G^2$ increases, i.e., more complex agents require more massive centers.

(b) The bandwidth $w$ decreases (Theorem~\ref{thm-freedom-bandwidth}c): more complex agents have narrower freedom bands.

(c) The coexistence attractor $r^*$ is unchanged: the optimal coupling structure is independent of agent complexity.
\end{proposition}

\textbf{Physical interpretation:} A complex, high-maintenance entity (a
research university, a democratic institution, a highly differentiated
organism) requires a richer environment to survive and has less margin
for error in its coupling distance. Simple, low-maintenance entities (a
single-celled organism, a subsistence farmer) can survive in sparser
environments with wider latitude. This is an \textbf{inherent trade-off}
between complexity and freedom: more complex agents require more
resources and are therefore more constrained in their viable coupling
range.

\subsection{Inequality: Asymmetric Agents at the Same
Center}\label{inequality-asymmetric-agents-at-the-same-center}

Consider two agents, \(i\) and \(j\), with different boundary
integrities \(B_i > B_j\) (agent \(i\) is more complex) coupled to the
same HEC. Both share the same attractor \(r^* = 2\tau/G\) and the same
HEC mass \(\mathcal{M}\).

\begin{corollary}[Inequality of Freedom]
\label{cor-inequality-of-freedom}

If $B_i > B_j$ (agent $i$ has higher boundary cost), then:

(a) $w_i < w_j$: agent $i$ has a narrower Coexistence Band (less freedom).

(b) $\mathcal{M}_{\min,i} > \mathcal{M}_{\min,j}$: agent $i$ requires a more massive center to survive.

(c) For any $\mathcal{M}$ with $\mathcal{M}_{\min,j} < \mathcal{M} < \mathcal{M}_{\min,i}$, agent $j$ can survive near this HEC but agent $i$ cannot.
\end{corollary}

\textbf{Physical interpretation:} Inequality in complexity creates
inequality in freedom. Within the same economic or political system,
agents with higher maintenance costs (due to greater complexity,
specialization, or vulnerability) have narrower viable ranges. This
formalization captures why developing nations have less policy latitude
than developed ones, why small businesses are more fragile than large
ones, and why specialized organisms are more extinction-prone than
generalists.

\section{Multi-Center Dynamics}\label{multi-center-dynamics}

\subsection{The Multi-Center
Potential}\label{the-multi-center-potential}

In realistic environments, agents interact with multiple HECs
simultaneously: an individual may be employed by a corporation, reside
in a state, trade with international partners, and participate in online
communities. Let there be \(K\) HECs with value masses
\(\mathcal{M}_1, \ldots, \mathcal{M}_K\).

The agent's coupling state is now a vector
\(\mathbf{r} = (r_1, \ldots, r_K) \in (0, \infty)^K\), where \(r_k\) is
the coupling distance to HEC \(k\).

\begin{definition}[Multi-Center Coexistence Potential]
\label{def-multi-center-potential}

The \emph{multi-center coexistence potential} is:

$$V(\mathbf{r}) = \sum_{k=1}^{K} \left[\frac{\tau_k \mathcal{M}_k}{r_k^2} - \frac{G_k \mathcal{M}_k}{r_k}\right] + \gamma B_i$$

where $G_k, \tau_k > 0$ are the center-specific coupling coefficients. The multi-center bandwidth is the volume of the viable set $\mathcal{B} = \{\mathbf{r} : V(\mathbf{r}) < 0\}$.
\end{definition}

\begin{theorem}[Multi-Center Coexistence Attractor]
\label{thm-multi-center-attractor}

The multi-center coexistence potential $V(\mathbf{r})$ has a unique global minimum at:

$$r_k^* = \frac{2\tau_k}{G_k} \quad \text{for each } k = 1, \ldots, K$$

This point is globally attracting under the gradient dynamics $\dot{\mathbf{r}} = -\mu \nabla V(\mathbf{r})$. The minimum value is:

$$V(\mathbf{r}^*) = \gamma B_i - \sum_{k=1}^{K} \frac{G_k^2 \mathcal{M}_k}{4\tau_k}$$

The multi-center Coexistence Band exists if and only if:

$$\sum_{k=1}^{K} \frac{G_k^2 \mathcal{M}_k}{4\tau_k} > \gamma B_i$$
\end{theorem}

\begin{proof}
The multi-center potential separates across the $K$ coupling dimensions: $V(\mathbf{r}) = \sum_k V_k(r_k) + \gamma B_i$ where $V_k(r_k) = \tau_k \mathcal{M}_k / r_k^2 - G_k \mathcal{M}_k / r_k$. Each $V_k$ has a unique minimum at $r_k^* = 2\tau_k/G_k$ with value $-G_k^2\mathcal{M}_k/(4\tau_k)$ (by the single-center analysis). The global minimum is achieved component-wise. Global attraction follows from the Lyapunov function $\mathcal{W}(\mathbf{r}) = V(\mathbf{r}) - V(\mathbf{r}^*)$, whose time derivative satisfies $\dot{\mathcal{W}} = -\mu\|\nabla V\|^2 \leq 0$. The existence condition requires $V(\mathbf{r}^*) < 0$.
\end{proof}

\begin{corollary}[Diversification Benefit]
\label{cor-diversification-benefit}

An agent coupled to $K$ HECs may be viable even if no single HEC has sufficient mass: it is possible that $\mathcal{M}_k < \mathcal{M}_{\min,k} \triangleq 4\gamma B_i \tau_k / G_k^2$ for each $k$ individually (the center-specific viability threshold of Theorem~\ref{thm-existence-coexistence-band}), yet $\sum_k G_k^2\mathcal{M}_k/(4\tau_k) > \gamma B_i$. Diversification across centers expands freedom.
\end{corollary}

\textbf{Physical interpretation:} An individual who cannot survive on
one job (\(\mathcal{M}_1 < \mathcal{M}_{\min}\)) may thrive by combining
income streams. A small nation unable to survive in the orbit of one
superpower may thrive through diversified trade agreements. The
multi-center model formalizes the survival advantage of diversification.

\subsection{Center Competition for
Agents}\label{center-competition-for-agents}

From the center's perspective, agents in its Coexistence Band contribute
to the network (labor, innovation, information) that increases the
center's own value mass \(\mathcal{M}\). Centers therefore compete for
agents by adjusting their coupling parameters:

\begin{itemize}
\tightlist
\item
  \textbf{Increasing \(G\)} (better resource sharing): raises wages,
  improves public services.
\item
  \textbf{Decreasing \(\tau\)} (reducing dissolution pressure): granting
  more autonomy, reducing oppressive constraints.
\end{itemize}

Both adjustments widen the Coexistence Band and attract agents from
competing centers. This formalizes why nations and corporations compete
through quality of life, wages, and individual freedoms, as it is a
rational strategy in the multi-center attractor landscape.

\section{Connection to the Broader
Framework}\label{connection-to-the-broader-framework}

\subsection{Coupling the Dynamical Model to Game
Theory}\label{coupling-the-dynamical-model-to-game-theory}

The discount factor \(\delta\) from the repeated game (Definition 2
\citep{TC-V}) depends on the agent's planning horizon and environment
stability. Agents within the Coexistence Band of a stable HEC have:

\begin{itemize}
\tightlist
\item
  Higher effective \(\delta\) (more stable environment → longer
  effective planning horizon)
\item
  A discount factor \(\delta(r)\) that is maximized at the coexistence
  attractor and exceeds the (fixed) cooperation threshold \(\delta^*\)
  throughout a subband of \(\mathcal{B}\)
  (Proposition~\ref{prop-stability-cooperation-feedback}); the threshold
  itself is a constant of the repeated game, unchanged by the dynamics
\end{itemize}

This means agents at the coexistence attractor are more likely to
sustain cooperation, creating a \textbf{positive feedback loop}:
stability → cooperation → more stability.

\begin{proposition}[Stability–Cooperation Feedback]
\label{prop-stability-cooperation-feedback}

Let an agent's discount factor depend on its net energy surplus: $\delta(r) = 1 - \exp(-\beta \max(\Pi(r), 0))$ for some $\beta > 0$ (agents with more surplus plan further ahead). Then:

(a) $\delta(r)$ is maximized at $r^*$ (the coexistence attractor).

(b) $\delta(r^*)> \delta(r)$ for all $r \neq r^*$ in $\mathcal{B}$.

(c) If $\delta(r^*) > \delta^*$ (equivalently, $\beta \Pi(r^*) > -\ln(1-\delta^*)$), then $\delta(r) > \delta^*$ throughout a subband of $\mathcal{B}$ centered on $r^*$, where $\delta^*$ is the cooperation threshold from Theorem 2 \citep{TC-V}.
\end{proposition}

\begin{proof}
We prove each claim in turn.

(a)–(b): $\Pi(r)$ is maximized at $r^*$ (by the definition of the attractor). Since $\delta(r) = 1 - e^{-\beta \Pi(r)}$ is strictly increasing in $\Pi(r)$, $\delta$ inherits the maximum. (c): Given $\delta(r^*) > \delta^*$, the fact that $\delta(r) \to 0$ continuously as $r \to r_\pm$ (since $\Pi(r) \to 0$ at the band boundaries) together with the intermediate value theorem gives a subband where $\delta(r) > \delta^*$, centered on $r^*$.
\end{proof}

\subsection{Coupling to Accumulated
Negentropy}\label{coupling-to-accumulated-negentropy}

The value mass \(\mathcal{M}\) is precisely the accumulated negentropy
defined in Definition 6 \citep{TC-III}:

\[\mathcal{M} = \mathcal{N}(T) = \int_0^T \dot{W}_{\text{order}}(t) \, dt\]

This means the Coexistence Band width is determined by the historical
thermodynamic investment in the center. A center that has accumulated
more negentropy over time provides wider freedom to its agents. The
\textbf{Irrationality of Destruction} (Theorem 6 \citep{TC-III})
acquires a new dimension here: destroying a high-\(\mathcal{M}\) center
does not just waste accumulated value; it \textbf{collapses the
Coexistence Bands} of all agents affiliated with that center,
catastrophically reducing the freedom of an entire population.

\begin{corollary}[Cascade Collapse from Center Destruction]
\label{cor-cascade-collapse}

If a HEC with value mass $\mathcal{M}$ is destroyed ($\mathcal{M} \to 0$), all $N$ agents within its Coexistence Band simultaneously lose viability with respect to that center; agents coupled solely to it lose viability outright, while diversified agents may survive on their remaining centers (Corollary~\ref{cor-diversification-benefit}). The total "freedom destroyed" is:

$$\mathcal{F}_{\text{lost}} = \sum_{i=1}^{N} w_i(\mathcal{M})$$

where $w_i$ is agent $i$'s freedom bandwidth; for homogeneous agents ($B_i \equiv B$) this reduces to $N \cdot w(\mathcal{M})$.

This provides an additional argument for the Irrationality of Destruction (Theorem 6 \citep{TC-III}): the social cost of center destruction includes the collapse of the entire affiliated population's viable coupling space.
\end{corollary}

\subsection{Coupling to Constrained
Optimization}\label{coupling-to-constrained-optimization}

The coupling distance \(r\) can be mapped to the constraint structure of
the constrained optimization derivation. At smaller \(r\) (tighter
coupling), the HEC imposes more boundary constraints (Definition 3
\citep{TC-I}) on the agent's action space, extracting a higher shadow
price \(\mu_{Bk}^*\) (Theorem 2 \citep{TC-I}). Here \(\mu_{Bk}^*\)
(abbreviated \(\mu_k^*\) below) denotes the boundary-constraint shadow
price of that theorem --- not the mobility \(\mu_0\) of the gradient
dynamics --- and \(m\) is the number of boundary constraints the HEC
imposes. The energetic inner boundary \(r_-\) corresponds to the
constraint intensity at which the aggregate shadow price of all
HEC-imposed constraints exceeds the agent's utility gain from resource
access:

\[\sum_{k=1}^{m} \mu_k^*(r_-) = U_i(\mathbf{x}_i^*) - U_i(\mathbf{0})\]

At this point, the constraints consume all of the agent's surplus: its
optimization space is so restricted that cooperation yields zero net
benefit.

\section{Worked Examples}\label{worked-examples}

\subsection{Individual
vs.~Corporation}\label{individual-vs.-corporation}

\textbf{Setup:} An individual agent (worker) with boundary integrity
\(B = 1\) (normalized), entropy leakage \(\gamma = 0.1\), near a
corporation with value mass \(\mathcal{M} = 100\) (normalized energy
units). Coupling parameters: \(G = 1.0\) (resource access),
\(\tau = 0.5\) (dissolution coupling).

\textbf{Calculations:}

\begin{itemize}
\tightlist
\item
  Coexistence attractor: \(r^* = 2\tau/G = 1.0\)
\item
  Minimum viable mass:
  \(\mathcal{M}_{\min} = 4\gamma B \tau / G^2 = 4(0.1)(1)(0.5)/1 = 0.2\),
  so the corporation far exceeds the threshold.
\item
  Discriminant:
  \(\Delta = G^2\mathcal{M}^2 - 4\gamma B \tau \mathcal{M} = 100^2 - 4(0.1)(1)(0.5)(100) = 10000 - 20 = 9980\)
\item
  Inner boundary:
  \(r_- = (100 - \sqrt{9980})/(2 \cdot 0.1 \cdot 1) = (100 - 99.90)/0.2 \approx 0.50\)
\item
  Outer boundary: \(r_+ = (100 + 99.90)/0.2 \approx 999.5\)
\item
  Bandwidth: \(w \approx 999.0\), an extremely wide freedom band.
\item
  Net energy at attractor:
  \(\Pi(r^*) = G^2\mathcal{M}/(4\tau) - \gamma B = 100/(2) - 0.1 = 49.9\)
\end{itemize}

\textbf{Interpretation:} The corporation is so massive relative to the
individual that the freedom band is enormous: the worker has a wide
range of viable engagement levels. The risk is not starvation (the outer
boundary is far away) but over-tight coupling: below
\(r_- \approx 0.50\) the energy balance turns negative. Active boundary
erosion is governed by the separate dissolution threshold \(r_d\)
(Theorem~\ref{thm-irreversibility-dissolution}), which depends
additionally on the assimilation intensity \(\sigma\) and the repair
parameters \(\rho\), \(\kappa\), and need not coincide with \(r_-\).

\subsection{Citizen vs.~State}\label{citizen-vs.-state}

\textbf{Setup:} A citizen with \(B = 2\) (higher complexity, including
education, social connections), \(\gamma = 0.05\) (stable social
environment), near a state with \(\mathcal{M} = 500\). Coupling:
\(G = 1.5\), \(\tau = 1.0\) (states impose more assimilation pressure
than corporations).

\begin{itemize}
\tightlist
\item
  \(r^* = 2(1.0)/1.5 \approx 1.33\)
\item
  \(\mathcal{M}_{\min} = 4(0.05)(2)(1.0)/1.5^2 = 0.4/2.25 \approx 0.178\)
\item
  \(\Delta = (1.5)^2(500)^2 - 4(0.05)(2)(1.0)(500) = 562500 - 200 = 562300\)
\item
  \(r_- \approx (750 - 749.87)/0.2 \approx 0.67\)
\item
  \(r_+ \approx (750 + 749.87)/0.2 \approx 7499\)
\item
  \(w \approx 7498\)
\end{itemize}

\textbf{Comparison: Authoritarian vs.~Free State.}

An authoritarian state may have higher \(\tau\) (more assimilation
pressure): \(\tau_{\text{auth}} = 3.0\).

\begin{itemize}
\tightlist
\item
  \(r^*_{\text{auth}} = 2(3.0)/1.5 = 4.0\) (agents must keep greater
  distance)
\item
  \(\mathcal{M}_{\min,\text{auth}} = 4(0.05)(2)(3.0)/2.25 = 0.533\)
  (requires more mass to sustain citizens)
\item
  \(\Delta_{\text{auth}} = 562500 - 600 = 561900\)
\item
  \(w_{\text{auth}} \approx \sqrt{561900}/0.1 \approx 7496\)
\end{itemize}

The bandwidth is similar in this regime
(\(\mathcal{M} \gg \mathcal{M}_{\min}\)), but the attractor distance
\(r^* = 4.0\) vs.~\(1.33\) means citizens of the authoritarian state
must maintain 3× more distance to avoid dissolution: they are
structurally pushed further from the center's resources.

\subsection{Small Nation
vs.~Superpower}\label{small-nation-vs.-superpower}

\textbf{Setup:} A small nation (agent) with \(B = 50\),
\(\gamma = 0.02\), near a superpower with \(\mathcal{M} = 10{,}000\).
Coupling: \(G = 0.5\) (trade only), \(\tau = 2.0\).

\begin{itemize}
\tightlist
\item
  \(r^* = 2(2.0)/0.5 = 8.0\)
\item
  \(\mathcal{M}_{\min} = 4(0.02)(50)(2.0)/0.25 = 32\)
\item
  \(\Delta = 0.25(10^8) - 4(0.02)(50)(2.0)(10^4) = 25{,}000{,}000 - 80{,}000 = 24{,}920{,}000\)
\item
  \(w = \sqrt{24920000}/(0.02 \times 50) = 4992.0/1.0 \approx 4992\)
\end{itemize}

If the superpower increases assimilation pressure (\(\tau \to 5\)):

\begin{itemize}
\tightlist
\item
  \(r^* = 20\) (nation must keep much greater distance)
\item
  \(\mathcal{M}_{\min} = 80\) (still viable, but threshold quintupled)
\item
  Bandwidth shrinks modestly but the viable zone shifts outward.
\end{itemize}

\textbf{Interpretation:} Small nations maintain sovereignty by keeping
sufficient coupling distance from superpowers, with trade and
cooperation at arm's length rather than full integration. When a
superpower increases its assimilation pressure (military threats,
economic coercion), the optimal distance increases and the energetic
inner boundary moves outward, shrinking the small nation's options.

\section{Discussion}\label{discussion}

\subsection{Summary of Results}\label{summary-of-results}

The paper derives a dynamical systems model of agent-center coupling
with three classes of results. First, the coexistence potential
(Definition~\ref{def-coexistence-potential}) constructed from a
short-range-dominant dissolution cost (\(1/r^2\)) and a long-range
resource inflow (\(1/r\)) possesses a unique global minimum, the
coexistence attractor (Definition~\ref{def-coexistence-attractor}),
which is globally asymptotically stable under gradient dynamics
(Theorem~\ref{thm-stability-coexistence-attractor}). The attractor
location \(r^* = 2\tau/G\) depends only on the dissolution-to-resource
coupling ratio, not on the center's mass or the agent's complexity.
Second, the stable coexistence band
(Definition~\ref{def-stable-coexistence-band}) exists if and only if the
center's value mass exceeds the threshold
\(\mathcal{M}_{\min} = 4\gamma B_i \tau / G^2\)
(Theorem~\ref{thm-existence-coexistence-band}), and the band's width,
the freedom bandwidth, admits a closed-form expression with monotonicity
in all model parameters (Theorem~\ref{thm-freedom-bandwidth}). Third,
the coupled agent-boundary system --- in which the boundary is an active
dissipative structure whose repair capacity and mobility scale with its
own integrity, and whose repair carries the multiplier \(\kappa > 1\) of
\citep{TC-IV} --- exhibits a basin of no return
(Theorem~\ref{thm-basin-of-no-return}): once boundary integrity falls
below a critical level \(B_c\), repair and mobility decay together and
the boundary reaches its collapse floor \(B_{\min}\) in finite time,
irreversibly. Dissolution --- sustained coupling below the dissolution
threshold \(r_d\), a threshold distinct from the energetic inner
boundary \(r_-\) (Theorem~\ref{thm-irreversibility-dissolution}) --- and
starvation --- sustained energy deficit outside the band
(Theorem~\ref{thm-starvation-spiral}) --- are the two routes into the
basin; transient excursions above the basin remain recoverable. The
multi-center extension (Theorem~\ref{thm-multi-center-attractor}) shows
that the potential separates across coupling dimensions, yielding a
componentwise global attractor with a diversification benefit
(Corollary~\ref{cor-diversification-benefit}).

\subsection{Relation to Potential Games and Gradient
Dynamics}\label{relation-to-potential-games-and-gradient-dynamics}

The agent-center model belongs to the class of gradient dynamical
systems on \((0, \infty)\) and inherits the convergence guarantees of
that theory \citep{Strogatz2015, HirschSmaleDevaney2012}. The
coexistence potential is a strict Lyapunov function for the dynamics, so
global convergence to the unique minimum is assured without additional
convexity assumptions on the payoff landscape.

The two-body potential form (short-range repulsion from dissolution,
long-range attraction from resource access) is structurally analogous to
the Lennard-Jones potential in molecular physics \citep{Jones1924}, with
the \(1/r^{12}\)-\(1/r^6\) pair replaced by \(1/r^2\)-\(1/r\). The lower
exponents reflect the longer interaction range of socio-economic
coupling relative to molecular forces. The analogy is mathematical, not
physical: we do not claim that literal intermolecular forces operate at
institutional scales. What is shared is the mathematical structure of a
potential with competing short-range and long-range terms that creates a
bounded stable basin between two catastrophic regimes.

The gradient dynamics assumption (\(\dot{r} = -\mu V'(r)\)) models
myopic adjustment: agents move in the direction of locally increasing
net energy without strategic foresight. This is appropriate for modeling
boundedly rational or habitual behavior and is formally analogous to
gradient play in the potential games literature
\citep{MondererShapley1996, Sandholm2010}, though the analogy is only
formal here: because each agent's payoff depends solely on its own
coupling coordinate --- the multi-center potential separates
(Theorem~\ref{thm-multi-center-attractor}) and agent-agent interactions
are absent from the model --- the aggregate potential decomposes into a
sum of independent single-agent potentials. The multi-agent system is
thus at most a degenerate (fully separable) potential game: \(N\)
decoupled gradient flows rather than a game with the genuine
cross-player strategic coupling that the potential-games literature
addresses, and gradient convergence follows agent-by-agent from the
single-agent Lyapunov argument rather than from potential-game
structure.

The model's dynamical predictions complement the static equilibrium
analysis of \emph{Cooperative Equilibrium} \citep{TC-V}. Where that
paper proves cooperation is the unique efficient Nash equilibrium, the
present paper proves that the dynamical adjustment process converges to
a stable cooperative configuration. The stability-cooperation feedback
loop (Proposition~\ref{prop-stability-cooperation-feedback}) connects
the two results: agents at the dynamical attractor sustain higher
discount factors, which in turn satisfy the cooperation threshold.

\subsection{Testable Predictions}\label{testable-predictions}

The model generates falsifiable predictions:

{\def\LTcaptype{none} % do not increment counter
\begin{longtable}[]{@{}
  >{\raggedright\arraybackslash}p{(\linewidth - 4\tabcolsep) * \real{0.3333}}
  >{\raggedright\arraybackslash}p{(\linewidth - 4\tabcolsep) * \real{0.3333}}
  >{\raggedright\arraybackslash}p{(\linewidth - 4\tabcolsep) * \real{0.3333}}@{}}
\toprule\noalign{}
\begin{minipage}[b]{\linewidth}\raggedright
Prediction
\end{minipage} & \begin{minipage}[b]{\linewidth}\raggedright
Source
\end{minipage} & \begin{minipage}[b]{\linewidth}\raggedright
Falsified If\ldots{}
\end{minipage} \\
\midrule\noalign{}
\endhead
\bottomrule\noalign{}
\endlastfoot
Agents cluster near a characteristic coupling distance \(r^*\) that
depends on coupling structure, not center mass &
Theorem~\ref{thm-stability-coexistence-attractor},
Definition~\ref{def-coexistence-attractor} & Empirical firm-worker or
citizen-state engagement levels scale with institution size rather than
with the ratio of autonomy cost to resource access \\
Freedom bandwidth increases with center mass &
Theorem~\ref{thm-freedom-bandwidth}(b) & Citizens of wealthier states or
employees of larger firms consistently exhibit narrower (not wider)
ranges of viable engagement levels \\
Boundary collapse is irreversible once integrity falls below the basin
boundary \(B_c\), while transient deficit or penetration above the basin
is recoverable & Theorem~\ref{thm-basin-of-no-return},
Theorem~\ref{thm-irreversibility-dissolution},
Theorem~\ref{thm-starvation-spiral} & Severely depleted entities
routinely recover boundary integrity without external intervention, or
mildly depleted entities systematically fail to recover \\
Minimum viable center mass scales linearly with agent complexity \(B_i\)
& Theorem~\ref{thm-existence-coexistence-band} & High-complexity agents
(specialized organisms, research-intensive firms) thrive near small-mass
centers as readily as low-complexity agents \\
Diversified agents survive where single-source agents fail &
Corollary~\ref{cor-diversification-benefit} & Agents with distributed
coupling across multiple centers show no survival advantage over agents
coupled to a single center of equivalent total mass \\
Cascade collapse: center destruction simultaneously eliminates the
viability of all agents coupled solely to that center &
Corollary~\ref{cor-cascade-collapse} & Destruction of a major
institution leaves the viable operating range of affiliated agents
unaffected \\
\end{longtable}
}

\subsection{Applications}\label{applications}

The results apply wherever self-maintaining agents couple to
resource-providing centers.

\textbf{Organizational economics.} The freedom bandwidth provides a
quantitative measure of the autonomy range available to agents within an
institution. The comparative statics results
(Proposition~\ref{prop-value-mass-comparative-statics},
Proposition~\ref{prop-boundary-integrity-comparative-statics}) predict
that larger, wealthier organizations can sustain wider autonomy for
their members, and that higher-complexity agents (specialized
professionals, research divisions) have narrower viable coupling ranges.
The dissolution and starvation thresholds correspond to observable
organizational pathologies: assimilation of subsidiaries and
isolation-induced decline, respectively. The center-competition
mechanism (§\ref{center-competition-for-agents}) formalizes why
institutions compete by increasing \(G\) (offering better resources) or
decreasing \(\tau\) (reducing conformity pressure), matching the
quality-of-life and autonomy dimensions that drive labor mobility.

\textbf{Institutional fragility and cascade collapse.} The cascade
collapse corollary (Corollary~\ref{cor-cascade-collapse}) predicts that
destruction of a high-mass center simultaneously eliminates the
viability of all agents coupled solely to that center. This formalizes
the empirical pattern in which institutional collapse (state failure,
corporate bankruptcy, ecosystem destruction) has outsized consequences
relative to the institution's own resource value, because the center's
coexistence band supports an entire population. The total freedom
destroyed, \(\sum_i w_i(\mathcal{M})\) --- \(N \cdot w(\mathcal{M})\)
for homogeneous agents --- provides a computable measure of systemic
risk, complementing the accumulated negentropy analysis of
\citep{TC-III}.

\textbf{Ecological niche modeling.} The stable coexistence band is a
mathematical analog of the ecological niche: a bounded region of
parameter space within which a species maintains positive net energy
\citep{Tilman1982}. The dependence of bandwidth on center mass parallels
the empirical observation that productive habitats support wider niche
breadth. The irreversibility results formalize the extinction vortex
\citep{Gilpin1986}, the self-reinforcing decline of populations that
fall below viability thresholds. The multi-center diversification
benefit provides a formal basis for the resilience advantage of species
with broader resource bases.

\subsection{Limitations}\label{limitations}

The model makes several simplifying assumptions. The coupling distance
\(r\) is a scalar, whereas real agent-center coupling is
multi-dimensional (economic, informational, political). A vector-valued
extension \(\mathbf{r} \in \mathbb{R}^k\) would better represent this
structure but would complicate the potential landscape and potentially
introduce multiple local minima, depending on the cross-term structure.

The gradient dynamics assumption models myopic adjustment. Agents with
strategic foresight would solve an optimal control problem rather than
following the gradient, potentially producing different transient
behavior. However, since the coexistence attractor is the unique global
minimum, the asymptotic destination is the same under any reasonable
adjustment rule.

The model parameters (\(G\), \(\tau\), \(\gamma\), \(\sigma\), \(\rho\),
\(\kappa\)) are treated as exogenous constants. In practice, all of
these are endogenous: centers adjust their coupling parameters in
response to agent behavior (§\ref{center-competition-for-agents}), and
agents' boundary integrity co-evolves with their coupling state
(§\ref{boundary-dynamics-and-irreversibility}). The coupled
agent-boundary system (§\ref{the-coupled-agent-boundary-system})
partially addresses this for \(B_i\) but leaves \(G\) and \(\tau\)
fixed.

Agent-agent interactions within a center's coexistence band are absent
from the model. In reality, agents coupled to the same center influence
one another through competition for the center's finite resources and
through network effects. Incorporating these interactions would couple
the \(N\) agents' dynamics and could produce coordination phenomena
(clustering, stratification) not captured by the single-agent analysis.

The separability of the multi-center potential
(Theorem~\ref{thm-multi-center-attractor}) assumes no center-center
interaction effects: coupling more tightly to center \(k\) does not
alter the parameters of center \(k'\). In practice, centers may compete
(joining one firm may constrain one's relationship with a rival) or
complement (a university affiliation may enhance corporate coupling).
Cross-coupling terms would break separability and could create multiple
local minima.

\textbf{Scope.} The results apply to any system of self-maintaining
agents coupled to resource-providing centers in a finite-resource
environment. The mathematical structure (potential well with short-range
repulsion and long-range attraction) is generic; the specific functional
forms (\(1/r^2\), \(1/r\)) are the simplest producing the correct
qualitative behavior. Different exponent choices would shift the
quantitative predictions (attractor location, bandwidth formula) while
preserving the structural results (unique attractor, bounded coexistence
band, irreversible transitions).

\subsection{Future Work}\label{future-work}

Several extensions follow naturally from the present results.

\textbf{Vector-valued coupling.} Replacing the scalar \(r\) with a
vector \(\mathbf{r} \in \mathbb{R}^k\) representing distinct coupling
dimensions (economic, informational, political) would permit analysis of
dimensional trade-offs: an agent might be closely coupled economically
but loosely coupled politically. The potential landscape in higher
dimensions may exhibit saddle points and ridges absent from the 1D
analysis, and the gradient dynamics would produce richer trajectory
geometry.

\textbf{Stochastic perturbations.} Adding noise to the gradient dynamics
(\(\dot{r} = -\mu V'(r) + \sigma_n \xi(t)\) with \(\xi\) a white noise
process) would model exogenous shocks (economic crises, policy changes,
natural disasters) and permit analysis of mean first-passage times from
the attractor to the band boundaries. The depth of the potential well
(\(D = G^2\mathcal{M}/(4\tau)\)) relative to noise intensity would
determine the system's resilience to stochastic perturbation, connecting
to the Kramers escape rate in chemical physics.

\textbf{Endogenous parameter dynamics.} Modeling \(G\) and \(\tau\) as
functions of system state (e.g., \(G\) increases with the number of
agents in the coexistence band due to network effects, or \(\tau\)
increases with \(\mathcal{M}\) due to bureaucratic growth) would
introduce feedback loops that could create hysteresis, bifurcations, or
limit cycles in the agent-center dynamics.

\textbf{Agent-agent interactions.} Incorporating coupling between agents
within a center's coexistence band, either through shared resource
depletion or through cooperative surplus generation, would extend the
single-agent model to a multi-agent dynamical system. Such an extension
could connect to evolutionary game dynamics on networks and to the
mean-field theory of interacting particle systems.

\textbf{Empirical calibration.} The model parameters (\(G\), \(\tau\),
\(\gamma\), \(B_i\)) are in principle measurable from institutional
data: resource distributions, autonomy metrics, organizational turnover
rates, and agent-complexity proxies. Calibrating the model against
firm-worker, citizen-state, or ecological data would convert the
qualitative predictions into quantitative forecasts and provide
empirical tests of the bandwidth scaling laws.

\textbf{Coupled agent-resource dynamics on production networks.} The
present dynamical analysis treats resource stocks as parameters. A
natural extension would couple agent-state dynamics to resource-stock
dynamics through a flow equation in which each agent's output
replenishes downstream inputs. This would double the phase space and
admit new phenomena: production cycles (boom-bust oscillations), trophic
cascades (single-agent dissolution propagating through the network),
hysteresis (irreversibility of collapse when producer agents dissolve,
connecting to \citep{TC-III}), and spectral stability conditions
requiring the network's input-output matrix to have spectral radius
below unity. Freedom bandwidth would acquire a supply-chain component,
giving a formal measure of network fragility.

\section{Conclusion}\label{conclusion}

We have derived a dynamical systems model of agent-center coupling in
which stable coexistence is a global attractor, the viable coupling
region is bounded and measurable, and boundary transitions are
irreversible. The coexistence attractor \(r^* = 2\tau/G\) is globally
asymptotically stable
(Theorem~\ref{thm-stability-coexistence-attractor}), with a convergence
rate proportional to the center's value mass and the agent's mobility.
The stable coexistence band exists if and only if the center's
accumulated negentropy exceeds a threshold that scales with agent
complexity (Theorem~\ref{thm-existence-coexistence-band}), and its
width, the freedom bandwidth, increases with center mass and decreases
with dissolution pressure and entropy leakage
(Theorem~\ref{thm-freedom-bandwidth}). Under the coupled agent-boundary
dynamics, in which repair capacity and mobility scale with boundary
integrity itself, collapse is governed by a basin of no return
(Theorem~\ref{thm-basin-of-no-return}): dissolution
(Theorem~\ref{thm-irreversibility-dissolution}) and starvation
(Theorem~\ref{thm-starvation-spiral}) become irreversible once depletion
crosses the basin boundary \(B_c\), and then complete in finite time.
The multi-center extension (Theorem~\ref{thm-multi-center-attractor})
proves that the potential separates across coupling dimensions, yielding
componentwise convergence and a diversification benefit for agents
coupled to multiple centers.

These results connect to and extend the preceding papers in this series.
The value mass \(\mathcal{M}\) that determines the coexistence band is
the accumulated negentropy of \citep{TC-III}; the center destruction
whose cascade consequences are quantified here
(Corollary~\ref{cor-cascade-collapse}) illustrates a dynamical
consequence of the irreplaceability result proved there. The boundary
integrity \(B_i\) whose degradation drives the irreversibility results
is maintained against the thermodynamic friction analyzed in
\citep{TC-IV}, and its regeneration is priced by the repair multiplier
\(\kappa > 1\) derived there. The stability-cooperation feedback
(Proposition~\ref{prop-stability-cooperation-feedback}) links the
dynamical attractor to the cooperation threshold of \citep{TC-V},
closing a loop between static equilibrium selection and dynamical
convergence. The constrained optimization framework of \citep{TC-I}
provides the shadow-price interpretation of the dissolution threshold
(§\ref{coupling-to-constrained-optimization}), and the information
integrity results of \citep{TC-II} underpin the informational dimension
of agent-center coupling. Together, the six papers provide a unified
account of the thermodynamics of multi-agent cooperation: constraint
necessity, information integrity, irreplaceability, friction costs,
equilibrium selection, and dynamical stability.

\addcontentsline{toc}{section}{References}
\begin{thebibliography}{14}
\providecommand{\natexlab}[1]{#1}
\providecommand{\url}[1]{\texttt{#1}}
\expandafter\ifx\csname urlstyle\endcsname\relax
  \providecommand{\doi}[1]{doi: #1}\else
  \providecommand{\doi}{doi: \begingroup \urlstyle{rm}\Url}\fi

\bibitem[Lostracco(2026{\natexlab{a}})]{TC-I}
Keith Lostracco.
\newblock Thermodynamics of cooperation: Necessary constraints,
  2026{\natexlab{a}}.
\newblock URL \url{https://doi.org/10.5281/zenodo.19635523}.
\newblock \doi{10.5281/zenodo.19635523}.

\bibitem[Lostracco(2026{\natexlab{b}})]{TC-II}
Keith Lostracco.
\newblock Thermodynamics of cooperation: Strategic entropy injection,
  2026{\natexlab{b}}.
\newblock URL \url{https://doi.org/10.5281/zenodo.19635814}.
\newblock \doi{10.5281/zenodo.19635814}.

\bibitem[Lostracco(2026{\natexlab{c}})]{TC-III}
Keith Lostracco.
\newblock Thermodynamics of cooperation: Accumulated negentropy,
  2026{\natexlab{c}}.
\newblock URL \url{https://doi.org/10.5281/zenodo.19635836}.
\newblock \doi{10.5281/zenodo.19635836}.

\bibitem[Lostracco(2026{\natexlab{d}})]{TC-IV}
Keith Lostracco.
\newblock Thermodynamics of cooperation: Thermodynamic friction,
  2026{\natexlab{d}}.
\newblock URL \url{https://doi.org/10.5281/zenodo.19635850}.
\newblock \doi{10.5281/zenodo.19635850}.

\bibitem[Lostracco(2026{\natexlab{e}})]{TC-V}
Keith Lostracco.
\newblock Thermodynamics of cooperation: Cooperative equilibrium,
  2026{\natexlab{e}}.
\newblock URL \url{https://doi.org/10.5281/zenodo.19635856}.
\newblock \doi{10.5281/zenodo.19635856}.

\bibitem[Strogatz(2015)]{Strogatz2015}
Steven~H. Strogatz.
\newblock \emph{Nonlinear Dynamics and Chaos: With Applications to Physics,
  Biology, Chemistry, and Engineering}.
\newblock CRC Press, 2nd edition, 2015.

\bibitem[Hirsch et~al.(2012)Hirsch, Smale, and Devaney]{HirschSmaleDevaney2012}
Morris~W. Hirsch, Stephen Smale, and Robert~L. Devaney.
\newblock \emph{Differential Equations, Dynamical Systems, and an Introduction
  to Chaos}.
\newblock Academic Press, 3rd edition, 2012.

\bibitem[Jones(1924)]{Jones1924}
J.~E. Jones.
\newblock On the determination of molecular fields. {II}. {From} the equation
  of state of a gas.
\newblock \emph{Proceedings of the Royal Society of London. Series A},
  106\penalty0 (738):\penalty0 463--477, 1924.
\newblock \doi{10.1098/rspa.1924.0082}.

\bibitem[Tirole(1988)]{Tirole1988}
Jean Tirole.
\newblock \emph{The Theory of Industrial Organization}.
\newblock MIT Press, 1988.

\bibitem[Acemoglu and Robinson(2012)]{AcemogluRobinson2012}
Daron Acemoglu and James~A. Robinson.
\newblock \emph{Why Nations Fail: The Origins of Power, Prosperity, and
  Poverty}.
\newblock Crown Publishers, 2012.

\bibitem[Tilman(1982)]{Tilman1982}
David Tilman.
\newblock \emph{Resource Competition and Community Structure}.
\newblock Princeton University Press, 1982.

\bibitem[Gilpin and Soul{\'e}(1986)]{Gilpin1986}
Michael~E. Gilpin and Michael~E. Soul{\'e}.
\newblock Minimum viable populations: Processes of species extinction.
\newblock In Michael~E. Soul{\'e}, editor, \emph{Conservation Biology: The
  Science of Scarcity and Diversity}, pages 19--34. Sinauer Associates, 1986.

\bibitem[Monderer and Shapley(1996)]{MondererShapley1996}
Dov Monderer and Lloyd Shapley.
\newblock Potential games.
\newblock \emph{Games and Economic Behavior}, 14:\penalty0 124--143, 05 1996.
\newblock \doi{10.1006/game.1996.0044}.

\bibitem[Sandholm(2010)]{Sandholm2010}
William~H. Sandholm.
\newblock \emph{Population Games and Evolutionary Dynamics}.
\newblock MIT Press, Cambridge, MA, 2010.
\newblock ISBN 9780262195874.

\end{thebibliography}
\end{document}
